\documentclass[../main.tex]{subfiles}
\begin{document}

\chapter{プロセスとプロセス管理}
\lessoninfo{
  video={P2L1},
  textbook={OSTEP \cite{ostep} ch.~4--7，13；Silberschatz ら \cite{silberschatz2018} ch.~3},
  keywords={プロセス，アドレス空間，仮想アドレス，ページテーブル，プロセス制御ブロック，コンテキストスイッチ，ホットキャッシュ，プロセス状態，fork，exec，CPUスケジューラ，タイムスライス，プロセス間通信，共有メモリ},
}

第1章では，オペレーティングシステムを，アプリケーションに代わって
ハードウェアを抽象化し調停する層として述べた．その最も基本的な抽象化が
プロセス，すなわち実行中のアプリケーションの表現である．この章では，
プロセスとは何か，どのような状態からなるか，オペレーティングシステムが
その状態をどのように表現するか，そして多数のプロセスを同時にどのように
管理するか---プロセスの間で CPU を切り替え，プロセスを生成し，
スケジュールし，互いに通信させる---を述べる．

\section{プログラムとプロセス}
\label{sec:l02-process}

ディスクやフラッシュメモリに保存されたアプリケーションは
\emph{プログラム}である．すなわち，命令と初期データからなるファイルで
あり，それ自体は何もしない静的な実体である．\source{P2L1，プロセスとは
何か；OSTEP \cite{ostep} ch.~4；Silberschatz ら \cite{silberschatz2018}
ch.~3．}%
プログラムが起動されると，オペレーティングシステムはそれをメモリへ
ロードして実行を開始し，プログラムは能動的な実体となる．

\begin{definition}[プロセス]
  \label{def:l02-process}
  \term[ぷろせす]{プロセス}[process]とは，実行中のプログラムの
  インスタンスである．すなわち，メモリにロードされた，実行中の
  プログラムの状態である．プロセスは\emph{タスク}あるいは\emph{ジョブ}%
  とも呼ばれる．
\end{definition}

同じプログラムを2回起動すれば2つのプロセスが生じる．両者は同じコードを
実行するが，それぞれが自分の状態---例えば異なる入力データや，プログラム
内の異なる位置---を持つ．プロセスは\emph{能動的な}アプリケーションの
実行状態を表すが，それはプロセスが常に CPU を使っていることを意味し
ない．ユーザの入力を待っているかもしれないし，自分の順番を待つ間に別の
プロセスが CPU 上で実行されているかもしれない．プロセスは，実行状態と
データのほかに，読み書きの対象となる I/O デバイスなどのハードウェアも
必要としうる．

\section{プロセスのアドレス空間}
\label{sec:l02-address-space}

プロセスは，実行中のアプリケーションのすべての状態---そのコード，
データ，および実行中に生成するすべての変数---をカプセル化する．%
\source{P2L1，プロセスはどのように見えるか，プロセスのアドレス空間；
OSTEP ch.~4，13．}%
この状態の各要素は，そのアドレスによって一意に識別できなければならない．

\begin{definition}[アドレス空間]
  プロセスの\term[あどれすくうかん]{アドレス空間}[address space]とは，
  オペレーティングシステムがプロセスの状態をカプセル化するために用いる
  抽象化であり，$V_0$ から $V_{\max}$ までのアドレスの範囲と，それらの
  アドレスに格納された状態からなる．これはプロセスのメモリ上の表現で
  ある．
\end{definition}

アドレス空間内の状態にはいくつかの種類があり，
\cref{fig:l02-address-space} のように配置される．
\begin{description}
  \item[テキストとデータ]
    \term[てきすとせぐめんと]{テキストセグメント}[text segment]は
    プログラムのコード，すなわちコンパイラが生成した命令を保持する．
    プログラムのロード時に初期化される大域変数などの静的データを保持
    するのが\term[でーたせぐめんと]{データセグメント}[data segment]で
    ある．どちらも静的な状態であり，プロセスが最初にロードされた時点
    から存在する．
  \item[ヒープ] \term[ひーぷ]{ヒープ}[heap]は，プロセスが実行中に動的に
    生成する状態，例えば \texttt{malloc} で確保したメモリを保持する．
    上位アドレスに向かって伸びる．
  \item[スタック] \term[すたっく]{スタック}[stack]はアドレス空間の
    動的な部分であり，実行中に後入れ先出しの順序で伸縮する．プロセスが
    手続きを呼び出すと，呼び出し元を再開するために必要な状態，すなわち
    戻り番地と退避されたレジスタの値がスタックに積まれ，呼び出された
    手続きの局所変数がそこに確保される．手続きから戻ると，それらは
    すべて取り除かれる．スタックはアドレス空間の上端から始まり，下位
    アドレスに向かって伸びる．
\end{description}

\begin{figure}[htbp]
  \centering
  % Layout of the address space of a process.
% Usage: % Layout of the address space of a process.
% Usage: % Layout of the address space of a process.
% Usage: \input{figures/l02-address-space}   (width ~112 mm)
\begin{tikzpicture}[x=1mm, y=1mm, font=\footnotesize\sffamily,
  >={Stealth[length=1.8mm]},
  ann/.style={font=\scriptsize\sffamily, anchor=west, align=left, text width=62mm},
  brk/.style={thin, ln-muted}]
  % regions
  \draw[fill=ln-box] (12,0)  rectangle (42,10) node[midway] {\iflangja{テキスト（コード）}{text (code)}};
  \draw[fill=ln-box] (12,10) rectangle (42,20) node[midway] {\iflangja{データ}{data}};
  \draw[fill=ln-box] (12,20) rectangle (42,30) node[midway] {\iflangja{ヒープ}{heap}};
  \draw[fill=white]  (12,30) rectangle (42,50);
  \draw[fill=ln-box] (12,50) rectangle (42,62) node[midway] {\iflangja{スタック}{stack}};
  \node[font=\scriptsize\sffamily, text=ln-muted] at (27,40) {\iflangja{未割り当て}{unallocated}};
  % growth directions
  \draw[->, thick, ln-accent] (19,30) -- (19,36);
  \draw[->, thick, ln-accent] (35,50) -- (35,44);
  % address labels
  \node[anchor=east] at (11,0.8)  {$V_0$};
  \node[anchor=east] at (11,61.2) {$V_{\max}$};
  % annotations
  \draw[brk] (44,0.5) -- (46,0.5) -- (46,19.5) -- (44,19.5);
  \node[ann] at (47,10) {\iflangja{静的な状態：プロセスのロード時から存在する}%
    {static state: present from the moment the process is loaded}};
  \draw[brk] (44,20.5) -- (46,20.5) -- (46,29.5) -- (44,29.5);
  \node[ann] at (47,25) {\iflangja{実行中に動的に確保される；上位アドレスへ伸びる}%
    {allocated dynamically during execution; grows toward higher addresses}};
  \draw[brk] (44,50.5) -- (46,50.5) -- (46,61.5) -- (44,61.5);
  \node[ann] at (47,56) {\iflangja{手続き呼び出しとともに後入れ先出しで伸縮する；下位アドレスへ伸びる}%
    {grows and shrinks last-in first-out with procedure calls; grows toward lower addresses}};
\end{tikzpicture}
   (width ~112 mm)
\begin{tikzpicture}[x=1mm, y=1mm, font=\footnotesize\sffamily,
  >={Stealth[length=1.8mm]},
  ann/.style={font=\scriptsize\sffamily, anchor=west, align=left, text width=62mm},
  brk/.style={thin, ln-muted}]
  % regions
  \draw[fill=ln-box] (12,0)  rectangle (42,10) node[midway] {\iflangja{テキスト（コード）}{text (code)}};
  \draw[fill=ln-box] (12,10) rectangle (42,20) node[midway] {\iflangja{データ}{data}};
  \draw[fill=ln-box] (12,20) rectangle (42,30) node[midway] {\iflangja{ヒープ}{heap}};
  \draw[fill=white]  (12,30) rectangle (42,50);
  \draw[fill=ln-box] (12,50) rectangle (42,62) node[midway] {\iflangja{スタック}{stack}};
  \node[font=\scriptsize\sffamily, text=ln-muted] at (27,40) {\iflangja{未割り当て}{unallocated}};
  % growth directions
  \draw[->, thick, ln-accent] (19,30) -- (19,36);
  \draw[->, thick, ln-accent] (35,50) -- (35,44);
  % address labels
  \node[anchor=east] at (11,0.8)  {$V_0$};
  \node[anchor=east] at (11,61.2) {$V_{\max}$};
  % annotations
  \draw[brk] (44,0.5) -- (46,0.5) -- (46,19.5) -- (44,19.5);
  \node[ann] at (47,10) {\iflangja{静的な状態：プロセスのロード時から存在する}%
    {static state: present from the moment the process is loaded}};
  \draw[brk] (44,20.5) -- (46,20.5) -- (46,29.5) -- (44,29.5);
  \node[ann] at (47,25) {\iflangja{実行中に動的に確保される；上位アドレスへ伸びる}%
    {allocated dynamically during execution; grows toward higher addresses}};
  \draw[brk] (44,50.5) -- (46,50.5) -- (46,61.5) -- (44,61.5);
  \node[ann] at (47,56) {\iflangja{手続き呼び出しとともに後入れ先出しで伸縮する；下位アドレスへ伸びる}%
    {grows and shrinks last-in first-out with procedure calls; grows toward lower addresses}};
\end{tikzpicture}
   (width ~112 mm)
\begin{tikzpicture}[x=1mm, y=1mm, font=\footnotesize\sffamily,
  >={Stealth[length=1.8mm]},
  ann/.style={font=\scriptsize\sffamily, anchor=west, align=left, text width=62mm},
  brk/.style={thin, ln-muted}]
  % regions
  \draw[fill=ln-box] (12,0)  rectangle (42,10) node[midway] {\iflangja{テキスト（コード）}{text (code)}};
  \draw[fill=ln-box] (12,10) rectangle (42,20) node[midway] {\iflangja{データ}{data}};
  \draw[fill=ln-box] (12,20) rectangle (42,30) node[midway] {\iflangja{ヒープ}{heap}};
  \draw[fill=white]  (12,30) rectangle (42,50);
  \draw[fill=ln-box] (12,50) rectangle (42,62) node[midway] {\iflangja{スタック}{stack}};
  \node[font=\scriptsize\sffamily, text=ln-muted] at (27,40) {\iflangja{未割り当て}{unallocated}};
  % growth directions
  \draw[->, thick, ln-accent] (19,30) -- (19,36);
  \draw[->, thick, ln-accent] (35,50) -- (35,44);
  % address labels
  \node[anchor=east] at (11,0.8)  {$V_0$};
  \node[anchor=east] at (11,61.2) {$V_{\max}$};
  % annotations
  \draw[brk] (44,0.5) -- (46,0.5) -- (46,19.5) -- (44,19.5);
  \node[ann] at (47,10) {\iflangja{静的な状態：プロセスのロード時から存在する}%
    {static state: present from the moment the process is loaded}};
  \draw[brk] (44,20.5) -- (46,20.5) -- (46,29.5) -- (44,29.5);
  \node[ann] at (47,25) {\iflangja{実行中に動的に確保される；上位アドレスへ伸びる}%
    {allocated dynamically during execution; grows toward higher addresses}};
  \draw[brk] (44,50.5) -- (46,50.5) -- (46,61.5) -- (44,61.5);
  \node[ann] at (47,56) {\iflangja{手続き呼び出しとともに後入れ先出しで伸縮する；下位アドレスへ伸びる}%
    {grows and shrinks last-in first-out with procedure calls; grows toward lower addresses}};
\end{tikzpicture}

  \caption{プロセスのアドレス空間．ヒープとスタックは互いに向かって
    伸びる．その間のアドレスは割り当てられていない．}
  \label{fig:l02-address-space}
\end{figure}

\subsection{仮想アドレスと物理アドレス}

プロセスが用いる $V_0$ から $V_{\max}$ までのアドレスは，マシンの物理
メモリ上の実際の位置に対応している必要はない．\source{P2L1，プロセスの
アドレス空間；OSTEP ch.~13．}

\begin{definition}[仮想アドレスと物理アドレス]
  プロセスのアドレス空間内のアドレスであり，プロセス自身が用いるものを
  \term[かそうあどれす]{仮想アドレス}[virtual address]と呼ぶ．これに
  対して，マシンの物理メモリ（DRAM）上の位置のアドレスが
  \term[ぶつりあどれす]{物理アドレス}[physical address]である．
\end{definition}

メモリ管理ハードウェアとオペレーティングシステムは，プロセスが用いる
すべての仮想アドレスを物理アドレスへ変換する．両者の対応付けは
\term[ぺーじてーぶる]{ページテーブル}[page table]に記録され，
オペレーティングシステムがこれをプロセスごとに保持する．2種類の
アドレスを分離することで，アドレス空間の配置が物理メモリの配置から
解放される．

\begin{example}
  あるプロセスが変数 \texttt{x} を仮想アドレス \texttt{0x00402A10} に
  格納しているとする．そのページテーブルはこのアドレスを物理アドレス
  \texttt{0x1F3A0A10} へ対応付けている．プロセスが \texttt{x} を
  読み書きするときに使うのは \texttt{0x00402A10} であり，
  \texttt{0x1F3A0A10} への変換はプロセスの関与なしに行われる．%
  \tip{変換はページ単位---x86 では \SI{4}{\kibi\byte}---で行われるので，
  アドレスの下位 12 ビットは変換を素通りする．ここでの2つのアドレスが
  ともに \texttt{A10} で終わるのはそのためである．}
\end{example}

\subsection{アドレス空間とメモリ管理}

この対応付けは，物理メモリの使い方についてオペレーティングシステムに
自由を与える．それが必要な理由は2つある．\source{P2L1，アドレス空間と
メモリ管理．}第一に，プロセスがアドレス空間の全体を使うことはまれで
あり，ヒープとスタックの間の領域のように，大きな部分はある時点では
割り当てられていない．第二に，割り当てられたすべての状態を収めるだけの
物理メモリがないかもしれない．32 ビットの仮想アドレスでは1つのプロセスが
$2^{32}$ バイト，すなわち \SI{4}{\gibi\byte} をアドレス指定でき，
複数のプロセスを合わせれば，必要なメモリはマシンが備える量を容易に
超える．\margin{x86-64 ではハードウェアが 48 ビットのアドレスを変換し，
Linux はその範囲の下半分，\SI{128}{\tebi\byte} を各プロセスに与える．
近年のプロセッサの5段ページングはアドレスを 57 ビットへ拡張する
（Linux の \emph{5-level paging} 文書）．}

したがってオペレーティングシステムは，どのアドレス空間のどの部分を物理
メモリのどこに置くかを動的に決める．一部は一時的にディスクへ移され，
プロセスが必要としたときにメモリへ戻される（スワッピング，第8章）．
オペレーティングシステムは，各プロセスについて，アドレス空間の各部分が
現在どこに---メモリ上かディスク上か---あるかを知っていなければならない．
これがページテーブルに保持される情報である．同じ情報によって，メモリ
アクセスの検証，すなわちプロセスが使おうとしているアドレスへの
アクセスが実際に許されているかの確認も行える．

プロセスはそれぞれ自分のアドレス空間を持ち，異なるプロセスのアドレス
空間は同じ仮想アドレスの範囲にわたる．したがって同時に実行される2つの
プロセスがまったく同じ仮想アドレスを使うことがあり，それぞれのページ
テーブルがそれを異なる物理的な位置へ対応付ける
（\cref{fig:l02-address-mapping}）．仮想アドレスは，それを使うプロセスと
合わせて初めて位置を特定する．

\begin{figure}[htbp]
  \centering
  % Two processes with the same virtual address range, mapped by their page
% tables to different physical frames (one region is on disk).
% Usage: % Two processes with the same virtual address range, mapped by their page
% tables to different physical frames (one region is on disk).
% Usage: % Two processes with the same virtual address range, mapped by their page
% tables to different physical frames (one region is on disk).
% Usage: \input{figures/l02-address-mapping}   (width ~112 mm)
\begin{tikzpicture}[x=1mm, y=1mm, font=\footnotesize\sffamily,
  >={Stealth[length=1.6mm]},
  pa/.style={fill=ln-accent!22},  % pages of P1
  pb/.style={fill=ln-source!18},  % pages of P2
  lab/.style={font=\scriptsize\sffamily},
  map/.style={->, thin},
  ttl/.style={font=\footnotesize\sffamily\bfseries, anchor=south}]
  % ---- virtual address spaces (same range for both processes)
  \foreach \x/\s/\n in {5/pa/P1, 91/pb/P2} {
    \draw[\s] (\x,0)  rectangle ++(14,8)  node[midway, lab] {\iflangja{テキスト}{text}};
    \draw[\s] (\x,8)  rectangle ++(14,8)  node[midway, lab] {\iflangja{データ}{data}};
    \draw[\s] (\x,16) rectangle ++(14,8)  node[midway, lab] {\iflangja{ヒープ}{heap}};
    \draw[fill=white] (\x,24) rectangle ++(14,14);
    \draw[\s] (\x,38) rectangle ++(14,10) node[midway, lab] {\iflangja{スタック}{stack}};
    \node[ttl] at (\x+7,49) {\n};
  }
  \node[lab, anchor=east] at (4.5,0.8)  {$V_0$};
  \node[lab, anchor=east] at (4.5,47.2) {$V_{\max}$};
  \node[lab, anchor=west] at (105.5,0.8)  {$V_0$};
  \node[lab, anchor=west] at (105.5,47.2) {$V_{\max}$};
  % ---- page tables
  \draw[fill=white] (25,0) rectangle (35,48)
    node[midway, rotate=90, lab] {\iflangja{P1 のページテーブル}{page table of P1}};
  \draw[fill=white] (75,0) rectangle (85,48)
    node[midway, rotate=90, lab] {\iflangja{P2 のページテーブル}{page table of P2}};
  % ---- physical memory and disk
  \node[ttl] at (55,49) {\iflangja{物理メモリ}{physical memory}};
  \foreach \i/\s in {0/pa, 1/pb, 2/pa, 3/pb, 4/pa, 5/pb, 6/pa} {
    \draw[\s] (47,13+5*\i) rectangle ++(16,5);
  }
  \draw[fill=ln-box, rounded corners=2pt] (47,0) rectangle (63,8)
    node[midway, lab] {\iflangja{ディスク}{disk}};
  % ---- P1 mappings: text->f0, data->f2, heap->f4, stack->f6
  \foreach \y/\f in {4/15.5, 12/25.5, 20/35.5, 43/45.5} {
    \draw (19,\y) -- (25,\y);
    \draw[map] (35,\y) -- (47,\f);
  }
  % ---- P2 mappings: text->disk, data->f1, heap->f3, stack->f5
  \foreach \y/\f in {4/4, 12/20.5, 20/30.5, 43/40.5} {
    \draw (91,\y) -- (85,\y);
    \draw[map] (75,\y) -- (63,\f);
  }
\end{tikzpicture}
   (width ~112 mm)
\begin{tikzpicture}[x=1mm, y=1mm, font=\footnotesize\sffamily,
  >={Stealth[length=1.6mm]},
  pa/.style={fill=ln-accent!22},  % pages of P1
  pb/.style={fill=ln-source!18},  % pages of P2
  lab/.style={font=\scriptsize\sffamily},
  map/.style={->, thin},
  ttl/.style={font=\footnotesize\sffamily\bfseries, anchor=south}]
  % ---- virtual address spaces (same range for both processes)
  \foreach \x/\s/\n in {5/pa/P1, 91/pb/P2} {
    \draw[\s] (\x,0)  rectangle ++(14,8)  node[midway, lab] {\iflangja{テキスト}{text}};
    \draw[\s] (\x,8)  rectangle ++(14,8)  node[midway, lab] {\iflangja{データ}{data}};
    \draw[\s] (\x,16) rectangle ++(14,8)  node[midway, lab] {\iflangja{ヒープ}{heap}};
    \draw[fill=white] (\x,24) rectangle ++(14,14);
    \draw[\s] (\x,38) rectangle ++(14,10) node[midway, lab] {\iflangja{スタック}{stack}};
    \node[ttl] at (\x+7,49) {\n};
  }
  \node[lab, anchor=east] at (4.5,0.8)  {$V_0$};
  \node[lab, anchor=east] at (4.5,47.2) {$V_{\max}$};
  \node[lab, anchor=west] at (105.5,0.8)  {$V_0$};
  \node[lab, anchor=west] at (105.5,47.2) {$V_{\max}$};
  % ---- page tables
  \draw[fill=white] (25,0) rectangle (35,48)
    node[midway, rotate=90, lab] {\iflangja{P1 のページテーブル}{page table of P1}};
  \draw[fill=white] (75,0) rectangle (85,48)
    node[midway, rotate=90, lab] {\iflangja{P2 のページテーブル}{page table of P2}};
  % ---- physical memory and disk
  \node[ttl] at (55,49) {\iflangja{物理メモリ}{physical memory}};
  \foreach \i/\s in {0/pa, 1/pb, 2/pa, 3/pb, 4/pa, 5/pb, 6/pa} {
    \draw[\s] (47,13+5*\i) rectangle ++(16,5);
  }
  \draw[fill=ln-box, rounded corners=2pt] (47,0) rectangle (63,8)
    node[midway, lab] {\iflangja{ディスク}{disk}};
  % ---- P1 mappings: text->f0, data->f2, heap->f4, stack->f6
  \foreach \y/\f in {4/15.5, 12/25.5, 20/35.5, 43/45.5} {
    \draw (19,\y) -- (25,\y);
    \draw[map] (35,\y) -- (47,\f);
  }
  % ---- P2 mappings: text->disk, data->f1, heap->f3, stack->f5
  \foreach \y/\f in {4/4, 12/20.5, 20/30.5, 43/40.5} {
    \draw (91,\y) -- (85,\y);
    \draw[map] (75,\y) -- (63,\f);
  }
\end{tikzpicture}
   (width ~112 mm)
\begin{tikzpicture}[x=1mm, y=1mm, font=\footnotesize\sffamily,
  >={Stealth[length=1.6mm]},
  pa/.style={fill=ln-accent!22},  % pages of P1
  pb/.style={fill=ln-source!18},  % pages of P2
  lab/.style={font=\scriptsize\sffamily},
  map/.style={->, thin},
  ttl/.style={font=\footnotesize\sffamily\bfseries, anchor=south}]
  % ---- virtual address spaces (same range for both processes)
  \foreach \x/\s/\n in {5/pa/P1, 91/pb/P2} {
    \draw[\s] (\x,0)  rectangle ++(14,8)  node[midway, lab] {\iflangja{テキスト}{text}};
    \draw[\s] (\x,8)  rectangle ++(14,8)  node[midway, lab] {\iflangja{データ}{data}};
    \draw[\s] (\x,16) rectangle ++(14,8)  node[midway, lab] {\iflangja{ヒープ}{heap}};
    \draw[fill=white] (\x,24) rectangle ++(14,14);
    \draw[\s] (\x,38) rectangle ++(14,10) node[midway, lab] {\iflangja{スタック}{stack}};
    \node[ttl] at (\x+7,49) {\n};
  }
  \node[lab, anchor=east] at (4.5,0.8)  {$V_0$};
  \node[lab, anchor=east] at (4.5,47.2) {$V_{\max}$};
  \node[lab, anchor=west] at (105.5,0.8)  {$V_0$};
  \node[lab, anchor=west] at (105.5,47.2) {$V_{\max}$};
  % ---- page tables
  \draw[fill=white] (25,0) rectangle (35,48)
    node[midway, rotate=90, lab] {\iflangja{P1 のページテーブル}{page table of P1}};
  \draw[fill=white] (75,0) rectangle (85,48)
    node[midway, rotate=90, lab] {\iflangja{P2 のページテーブル}{page table of P2}};
  % ---- physical memory and disk
  \node[ttl] at (55,49) {\iflangja{物理メモリ}{physical memory}};
  \foreach \i/\s in {0/pa, 1/pb, 2/pa, 3/pb, 4/pa, 5/pb, 6/pa} {
    \draw[\s] (47,13+5*\i) rectangle ++(16,5);
  }
  \draw[fill=ln-box, rounded corners=2pt] (47,0) rectangle (63,8)
    node[midway, lab] {\iflangja{ディスク}{disk}};
  % ---- P1 mappings: text->f0, data->f2, heap->f4, stack->f6
  \foreach \y/\f in {4/15.5, 12/25.5, 20/35.5, 43/45.5} {
    \draw (19,\y) -- (25,\y);
    \draw[map] (35,\y) -- (47,\f);
  }
  % ---- P2 mappings: text->disk, data->f1, heap->f3, stack->f5
  \foreach \y/\f in {4/4, 12/20.5, 20/30.5, 43/40.5} {
    \draw (91,\y) -- (85,\y);
    \draw[map] (75,\y) -- (63,\f);
  }
\end{tikzpicture}

  \caption{同じ仮想アドレスの範囲を持つ2つのプロセス．それぞれの
    ページテーブルは，割り当てられた領域を異なる物理フレームへ対応付け
    ている．P2 のテキストは現在ディスク上にあり，割り当てられていない
    領域はどこにも対応付けられていない．}
  \label{fig:l02-address-mapping}
\end{figure}

\section{実行状態とプロセス制御ブロック}
\label{sec:l02-pcb}

オペレーティングシステムがプロセスを止めるときには，後でまったく同じ
地点からそのプロセスを再開できるように，プロセスが何をしていたかを記録
しなければならない．\source{P2L1，プロセスの実行状態，プロセス制御
ブロックとは何か；OSTEP ch.~4．}プログラムはバイナリ，すなわち命令の列
にコンパイルされるが，その命令が厳密に順番どおりに実行されるとは限ら
ない．ジャンプ，ループ，割り込みが順序を変える．CPU は，この列のどこに
プロセスがいるかを常に知っていなければならない．

\term[ぷろぐらむかうんた]{プログラムカウンタ}[program counter]%
\index{PC|see{プログラムカウンタ}}（PC）は，次に実行する命令のアドレスを
保持する．プロセスの実行中は，CPU のレジスタに保持される．CPU のその他の
レジスタは，データの値，アドレス，そして実行の順序に影響する状態情報を
保持する．\term[すたっくぽいんた]{スタックポインタ}[stack pointer]（SP）
はスタックの先頭のアドレスを保持する．これらの値が合わさってプロセスの
実行状態をなす．

\begin{definition}[プロセス制御ブロック]
  \term[ぷろせすせいぎょぶろっく]{プロセス制御ブロック}[process control block]%
  \index{PCB|see{プロセス制御ブロック}}（PCB）とは，オペレーティング
  システムがプロセスごとに保持するデータ構造である．プロセス状態，
  プロセス番号，プログラムカウンタとその他の CPU レジスタ，メモリの
  範囲と仮想から物理への対応付け，オープンしたファイルの一覧，優先度，
  シグナルマスク，CPU スケジューリング情報を含む．
\end{definition}

PCB の各フィールド（\cref{fig:l02-pcb}）は，異なる時点で維持される．%
\caution{Linux にプロセス専用の構造体はない．1つの
\texttt{struct task\_struct} が1つの\emph{スレッド}（第3章）を表し，
プロセスとはアドレス空間を共有するタスクの集まりである．}
\begin{itemize}
  \item PCB はプロセスの生成時に作られ，初期化される．例えばプログラム
    カウンタはプログラムの最初の命令に設定される．
  \item 一部のフィールドは，対応する状態が変化したときに更新される．
    プロセスがより多くのメモリを要求すると，オペレーティングシステムは
    それを割り当て，メモリの範囲と対応付けを更新する．
  \item その他のフィールドは，\margin{\texttt{/proc/PID/status} はこれらの
    フィールドの多く---状態，親，スレッド数，メモリの範囲，シグナル
    マスク---を表示する．x86-64 では各タスクが \SI{16}{\kibi\byte} の
    カーネルスタックも持つ（Linux の \emph{Kernel Stacks} 文書）．}%
    PCB 内で最新に保つにはあまりに頻繁に変化
    する．プログラムカウンタは命令ごとに変化するので，CPU が自分の
    レジスタに保持する．オペレーティングシステムは，CPU がプロセスの
    ために保持している値を集め，そのプロセスが CPU 上での実行を止める
    たびに PCB へ退避する．
\end{itemize}

\begin{figure}[htbp]
  \centering
  % Process control blocks of two processes and the CPU registers whose
% contents are saved into and restored from them.
% Usage: % Process control blocks of two processes and the CPU registers whose
% contents are saved into and restored from them.
% Usage: % Process control blocks of two processes and the CPU registers whose
% contents are saved into and restored from them.
% Usage: \input{figures/l02-pcb}   (width ~116 mm)
\begin{tikzpicture}[x=1mm, y=1mm, font=\scriptsize\sffamily,
  >={Stealth[length=1.6mm]},
  fld/.style={draw, minimum width=34mm, minimum height=4.5mm, inner sep=0pt, fill=white},
  cpuf/.style={fld, fill=ln-accent!22},
  reg/.style={draw, minimum width=24mm, minimum height=4.5mm, inner sep=0pt, fill=ln-accent!22},
  ttl/.style={font=\footnotesize\sffamily\bfseries, anchor=south}]
  \foreach \x/\n in {17/1, 99/2} {
    \node[ttl] at (\x,2.5) {\iflangja{P\n の PCB}{PCB of P\n}};
    \node[fld]  at (\x,0)     {\iflangja{プロセス状態}{process state}};
    \node[fld]  at (\x,-4.5)  {\iflangja{プロセス番号}{process number}};
    \node[cpuf] at (\x,-9)    {\iflangja{プログラムカウンタ}{program counter}};
    \node[cpuf] at (\x,-13.5) {\iflangja{レジスタ}{registers}};
    \node[fld]  at (\x,-18)   {\iflangja{メモリの範囲・対応付け}{memory limits / mappings}};
    \node[fld]  at (\x,-22.5) {\iflangja{オープンしたファイルの一覧}{list of open files}};
    \node[fld]  at (\x,-27)   {\iflangja{優先度}{priority}};
    \node[fld]  at (\x,-31.5) {\iflangja{シグナルマスク}{signal mask}};
    \node[fld]  at (\x,-36)   {\iflangja{CPU スケジューリング情報}{CPU scheduling info}};
    \node[fld]  at (\x,-40.5) {\ldots};
  }
  % CPU
  \draw[rounded corners=2pt, fill=ln-box] (46,-19) rectangle (70,2.5);
  \node[font=\footnotesize\sffamily\bfseries] at (58,-2) {CPU};
  \node[reg] at (58,-9)    {PC};
  \node[reg] at (58,-13.5) {\iflangja{レジスタ（SP を含む）}{registers (incl.\ SP)}};
  % save P1, restore P2
  \draw[->, thick] (46,-9) -- (34,-9);
  \draw[->, thick] (46,-13.5) -- (34,-13.5);
  \node[anchor=north, align=center] at (40,-15.5) {\iflangja{P1 の\\状態を\\退避}{save\\state\\of P1}};
  \draw[->, thick] (82,-9) -- (70,-9);
  \draw[->, thick] (82,-13.5) -- (70,-13.5);
  \node[anchor=north, align=center] at (76,-15.5) {\iflangja{P2 の\\状態を\\復元}{restore\\state\\of P2}};
\end{tikzpicture}
   (width ~116 mm)
\begin{tikzpicture}[x=1mm, y=1mm, font=\scriptsize\sffamily,
  >={Stealth[length=1.6mm]},
  fld/.style={draw, minimum width=34mm, minimum height=4.5mm, inner sep=0pt, fill=white},
  cpuf/.style={fld, fill=ln-accent!22},
  reg/.style={draw, minimum width=24mm, minimum height=4.5mm, inner sep=0pt, fill=ln-accent!22},
  ttl/.style={font=\footnotesize\sffamily\bfseries, anchor=south}]
  \foreach \x/\n in {17/1, 99/2} {
    \node[ttl] at (\x,2.5) {\iflangja{P\n の PCB}{PCB of P\n}};
    \node[fld]  at (\x,0)     {\iflangja{プロセス状態}{process state}};
    \node[fld]  at (\x,-4.5)  {\iflangja{プロセス番号}{process number}};
    \node[cpuf] at (\x,-9)    {\iflangja{プログラムカウンタ}{program counter}};
    \node[cpuf] at (\x,-13.5) {\iflangja{レジスタ}{registers}};
    \node[fld]  at (\x,-18)   {\iflangja{メモリの範囲・対応付け}{memory limits / mappings}};
    \node[fld]  at (\x,-22.5) {\iflangja{オープンしたファイルの一覧}{list of open files}};
    \node[fld]  at (\x,-27)   {\iflangja{優先度}{priority}};
    \node[fld]  at (\x,-31.5) {\iflangja{シグナルマスク}{signal mask}};
    \node[fld]  at (\x,-36)   {\iflangja{CPU スケジューリング情報}{CPU scheduling info}};
    \node[fld]  at (\x,-40.5) {\ldots};
  }
  % CPU
  \draw[rounded corners=2pt, fill=ln-box] (46,-19) rectangle (70,2.5);
  \node[font=\footnotesize\sffamily\bfseries] at (58,-2) {CPU};
  \node[reg] at (58,-9)    {PC};
  \node[reg] at (58,-13.5) {\iflangja{レジスタ（SP を含む）}{registers (incl.\ SP)}};
  % save P1, restore P2
  \draw[->, thick] (46,-9) -- (34,-9);
  \draw[->, thick] (46,-13.5) -- (34,-13.5);
  \node[anchor=north, align=center] at (40,-15.5) {\iflangja{P1 の\\状態を\\退避}{save\\state\\of P1}};
  \draw[->, thick] (82,-9) -- (70,-9);
  \draw[->, thick] (82,-13.5) -- (70,-13.5);
  \node[anchor=north, align=center] at (76,-15.5) {\iflangja{P2 の\\状態を\\復元}{restore\\state\\of P2}};
\end{tikzpicture}
   (width ~116 mm)
\begin{tikzpicture}[x=1mm, y=1mm, font=\scriptsize\sffamily,
  >={Stealth[length=1.6mm]},
  fld/.style={draw, minimum width=34mm, minimum height=4.5mm, inner sep=0pt, fill=white},
  cpuf/.style={fld, fill=ln-accent!22},
  reg/.style={draw, minimum width=24mm, minimum height=4.5mm, inner sep=0pt, fill=ln-accent!22},
  ttl/.style={font=\footnotesize\sffamily\bfseries, anchor=south}]
  \foreach \x/\n in {17/1, 99/2} {
    \node[ttl] at (\x,2.5) {\iflangja{P\n の PCB}{PCB of P\n}};
    \node[fld]  at (\x,0)     {\iflangja{プロセス状態}{process state}};
    \node[fld]  at (\x,-4.5)  {\iflangja{プロセス番号}{process number}};
    \node[cpuf] at (\x,-9)    {\iflangja{プログラムカウンタ}{program counter}};
    \node[cpuf] at (\x,-13.5) {\iflangja{レジスタ}{registers}};
    \node[fld]  at (\x,-18)   {\iflangja{メモリの範囲・対応付け}{memory limits / mappings}};
    \node[fld]  at (\x,-22.5) {\iflangja{オープンしたファイルの一覧}{list of open files}};
    \node[fld]  at (\x,-27)   {\iflangja{優先度}{priority}};
    \node[fld]  at (\x,-31.5) {\iflangja{シグナルマスク}{signal mask}};
    \node[fld]  at (\x,-36)   {\iflangja{CPU スケジューリング情報}{CPU scheduling info}};
    \node[fld]  at (\x,-40.5) {\ldots};
  }
  % CPU
  \draw[rounded corners=2pt, fill=ln-box] (46,-19) rectangle (70,2.5);
  \node[font=\footnotesize\sffamily\bfseries] at (58,-2) {CPU};
  \node[reg] at (58,-9)    {PC};
  \node[reg] at (58,-13.5) {\iflangja{レジスタ（SP を含む）}{registers (incl.\ SP)}};
  % save P1, restore P2
  \draw[->, thick] (46,-9) -- (34,-9);
  \draw[->, thick] (46,-13.5) -- (34,-13.5);
  \node[anchor=north, align=center] at (40,-15.5) {\iflangja{P1 の\\状態を\\退避}{save\\state\\of P1}};
  \draw[->, thick] (82,-9) -- (70,-9);
  \draw[->, thick] (82,-13.5) -- (70,-13.5);
  \node[anchor=north, align=center] at (76,-15.5) {\iflangja{P2 の\\状態を\\復元}{restore\\state\\of P2}};
\end{tikzpicture}

  \caption{P1 と P2 のプロセス制御ブロック．網掛けのフィールドは，
    プロセスの実行中に CPU がレジスタに保持する値であり，プロセスが
    実行を止めるときに PCB へ退避され，再び実行されるときに PCB から
    復元される．}
  \label{fig:l02-pcb}
\end{figure}

\section{コンテキストスイッチ}
\label{sec:l02-context-switch}

2つのプロセス P1 と P2 があり，その PCB がオペレーティングシステムの
メモリに格納されているとする．\source{P2L1，PCB はどのように使われるか，
コンテキストスイッチ；OSTEP ch.~6．}P1 の実行中，CPU のレジスタは P1 の
実行状態を保持している．オペレーティングシステムが P1 に割り込むと，
レジスタの内容を P1 の PCB へ退避し，P2 の状態を P2 の PCB から
レジスタへ復元する．そして P2 が実行され，P1 は休止する．後に
オペレーティングシステムが P2 に割り込むと，P2 の状態をその PCB へ
退避し，P1 の状態を復元する．P1 は，止められたまさにその地点から再開
する（\cref{fig:l02-context-switch}）．

\begin{definition}[コンテキストスイッチ]
  \term[こんてきすとすいっち]{コンテキストスイッチ}[context switch]とは，
  オペレーティングシステムが CPU をあるプロセスのコンテキストから別の
  プロセスのコンテキストへ切り替えるために用いる機構である．
\end{definition}

\begin{figure}[htbp]
  \centering
  % Context switches between two processes over time.
% Usage: % Context switches between two processes over time.
% Usage: % Context switches between two processes over time.
% Usage: \input{figures/l02-context-switch}   (width ~116 mm)
\begin{tikzpicture}[x=1mm, y=1mm, font=\footnotesize\sffamily,
  >={Stealth[length=1.6mm]},
  ex/.style={draw, fill=ln-accent!22, minimum height=7mm, inner sep=0pt},
  os/.style={draw, fill=ln-box, minimum height=7mm, inner sep=1pt, font=\scriptsize\sffamily, align=center},
  idle/.style={densely dashed, ln-rule, thick},
  lab/.style={font=\scriptsize\sffamily, text=ln-muted}]
  % row labels
  \node[anchor=east] at (11,24) {P1};
  \node[anchor=east] at (11,12) {OS};
  \node[anchor=east] at (11,0)  {P2};
  % P1 executing
  \node[ex, minimum width=24mm] at (24,24) {\iflangja{実行中}{executing}};
  \draw[idle] (36,24) -- (100,24);
  \node[ex, minimum width=16mm] at (108,24) {\iflangja{実行中}{executing}};
  % P2
  \draw[idle] (12,0) -- (56,0);
  \node[ex, minimum width=24mm] at (68,0) {\iflangja{実行中}{executing}};
  \draw[idle] (80,0) -- (116,0);
  % OS work
  \node[os, minimum width=20mm] (s1) at (46,12)
    {\iflangja{PCB\textsubscript{1} へ退避\\PCB\textsubscript{2} から復元}{save into PCB\textsubscript{1}\\restore PCB\textsubscript{2}}};
  \node[os, minimum width=20mm] (s2) at (90,12)
    {\iflangja{PCB\textsubscript{2} へ退避\\PCB\textsubscript{1} から復元}{save into PCB\textsubscript{2}\\restore PCB\textsubscript{1}}};
  % transfers of control
  \draw[->] (36,20.5) -- (36,15.5);
  \draw[->] (56,8.5)  -- (56,3.5);
  \draw[->] (80,3.5)  -- (80,8.5);
  \draw[->] (100,15.5) -- (100,20.5);
  \node[lab, anchor=east] at (35.5,18) {\iflangja{割り込み}{interrupt}};
  \node[lab, anchor=east] at (79.5,6)  {\iflangja{割り込み}{interrupt}};
  % context switch marks and time axis
  \node[lab, anchor=north] at (46,-5) {\iflangja{コンテキストスイッチ}{context switch}};
  \node[lab, anchor=north] at (90,-5) {\iflangja{コンテキストスイッチ}{context switch}};
  \draw[->, ln-muted] (12,-11) -- (116,-11) node[lab, anchor=north east] at (116,-11.5) {\iflangja{時間}{time}};
\end{tikzpicture}
   (width ~116 mm)
\begin{tikzpicture}[x=1mm, y=1mm, font=\footnotesize\sffamily,
  >={Stealth[length=1.6mm]},
  ex/.style={draw, fill=ln-accent!22, minimum height=7mm, inner sep=0pt},
  os/.style={draw, fill=ln-box, minimum height=7mm, inner sep=1pt, font=\scriptsize\sffamily, align=center},
  idle/.style={densely dashed, ln-rule, thick},
  lab/.style={font=\scriptsize\sffamily, text=ln-muted}]
  % row labels
  \node[anchor=east] at (11,24) {P1};
  \node[anchor=east] at (11,12) {OS};
  \node[anchor=east] at (11,0)  {P2};
  % P1 executing
  \node[ex, minimum width=24mm] at (24,24) {\iflangja{実行中}{executing}};
  \draw[idle] (36,24) -- (100,24);
  \node[ex, minimum width=16mm] at (108,24) {\iflangja{実行中}{executing}};
  % P2
  \draw[idle] (12,0) -- (56,0);
  \node[ex, minimum width=24mm] at (68,0) {\iflangja{実行中}{executing}};
  \draw[idle] (80,0) -- (116,0);
  % OS work
  \node[os, minimum width=20mm] (s1) at (46,12)
    {\iflangja{PCB\textsubscript{1} へ退避\\PCB\textsubscript{2} から復元}{save into PCB\textsubscript{1}\\restore PCB\textsubscript{2}}};
  \node[os, minimum width=20mm] (s2) at (90,12)
    {\iflangja{PCB\textsubscript{2} へ退避\\PCB\textsubscript{1} から復元}{save into PCB\textsubscript{2}\\restore PCB\textsubscript{1}}};
  % transfers of control
  \draw[->] (36,20.5) -- (36,15.5);
  \draw[->] (56,8.5)  -- (56,3.5);
  \draw[->] (80,3.5)  -- (80,8.5);
  \draw[->] (100,15.5) -- (100,20.5);
  \node[lab, anchor=east] at (35.5,18) {\iflangja{割り込み}{interrupt}};
  \node[lab, anchor=east] at (79.5,6)  {\iflangja{割り込み}{interrupt}};
  % context switch marks and time axis
  \node[lab, anchor=north] at (46,-5) {\iflangja{コンテキストスイッチ}{context switch}};
  \node[lab, anchor=north] at (90,-5) {\iflangja{コンテキストスイッチ}{context switch}};
  \draw[->, ln-muted] (12,-11) -- (116,-11) node[lab, anchor=north east] at (116,-11.5) {\iflangja{時間}{time}};
\end{tikzpicture}
   (width ~116 mm)
\begin{tikzpicture}[x=1mm, y=1mm, font=\footnotesize\sffamily,
  >={Stealth[length=1.6mm]},
  ex/.style={draw, fill=ln-accent!22, minimum height=7mm, inner sep=0pt},
  os/.style={draw, fill=ln-box, minimum height=7mm, inner sep=1pt, font=\scriptsize\sffamily, align=center},
  idle/.style={densely dashed, ln-rule, thick},
  lab/.style={font=\scriptsize\sffamily, text=ln-muted}]
  % row labels
  \node[anchor=east] at (11,24) {P1};
  \node[anchor=east] at (11,12) {OS};
  \node[anchor=east] at (11,0)  {P2};
  % P1 executing
  \node[ex, minimum width=24mm] at (24,24) {\iflangja{実行中}{executing}};
  \draw[idle] (36,24) -- (100,24);
  \node[ex, minimum width=16mm] at (108,24) {\iflangja{実行中}{executing}};
  % P2
  \draw[idle] (12,0) -- (56,0);
  \node[ex, minimum width=24mm] at (68,0) {\iflangja{実行中}{executing}};
  \draw[idle] (80,0) -- (116,0);
  % OS work
  \node[os, minimum width=20mm] (s1) at (46,12)
    {\iflangja{PCB\textsubscript{1} へ退避\\PCB\textsubscript{2} から復元}{save into PCB\textsubscript{1}\\restore PCB\textsubscript{2}}};
  \node[os, minimum width=20mm] (s2) at (90,12)
    {\iflangja{PCB\textsubscript{2} へ退避\\PCB\textsubscript{1} から復元}{save into PCB\textsubscript{2}\\restore PCB\textsubscript{1}}};
  % transfers of control
  \draw[->] (36,20.5) -- (36,15.5);
  \draw[->] (56,8.5)  -- (56,3.5);
  \draw[->] (80,3.5)  -- (80,8.5);
  \draw[->] (100,15.5) -- (100,20.5);
  \node[lab, anchor=east] at (35.5,18) {\iflangja{割り込み}{interrupt}};
  \node[lab, anchor=east] at (79.5,6)  {\iflangja{割り込み}{interrupt}};
  % context switch marks and time axis
  \node[lab, anchor=north] at (46,-5) {\iflangja{コンテキストスイッチ}{context switch}};
  \node[lab, anchor=north] at (90,-5) {\iflangja{コンテキストスイッチ}{context switch}};
  \draw[->, ln-muted] (12,-11) -- (116,-11) node[lab, anchor=north east] at (116,-11.5) {\iflangja{時間}{time}};
\end{tikzpicture}

  \caption{P1 と P2 の間の2回のコンテキストスイッチ．各プロセスは他方が
    実行される間は休止（破線）しており，その状態は PCB の中で待って
    いる．}
  \label{fig:l02-context-switch}
\end{figure}

コンテキストスイッチが高価である理由は2つある．
\begin{description}
  \item[直接のコスト] 一方のプロセスの状態を退避し他方の状態を復元する
    ロード命令とストア命令の実行に費やされる CPU サイクル．
  \item[間接のコスト] プロセッサキャッシュへの影響．キャッシュアクセス
    は数サイクル程度で済むのに対し，主記憶へのアクセスは数百サイクル
    程度を要する．\margin{Intel Skylake では L1 ヒットが 4--5 サイクル，
    L2 が約 12 サイクル，L3 が約 44 サイクル，DRAM はさらに
    \SI{50}{\nano\second} 程度（7-cpu.com）．}
\end{description}

プロセスの実行中は，そのプロセスが使うデータの多くがプロセッサ
キャッシュに存在する．そのようなキャッシュは
\term[ほっときゃっしゅ]{ホットキャッシュ}[hot cache]であり，ほとんどの
アクセスがキャッシュから供給され，プロセスは最良の性能を発揮する．
コンテキストスイッチの後，次のプロセスがこのデータの多くを自分の
データで置き換える．最初のプロセスが再開したときに見出すのは
\term[こーるどきゃっしゅ]{コールドキャッシュ}[cold cache]である．
自分のデータはもはやキャッシュになく，キャッシュが再び満たされるまで，
そのアクセスは主記憶から供給されるほかないキャッシュミスとなる．
したがってオペレーティングシステムは，コンテキストを切り替える頻度を
抑えるべきである．

\begin{example}
  レジスタの退避と復元に 1\,000 サイクルかかり，アクセスのコストは
  データがキャッシュにあるとき 4 サイクル，主記憶から読まなければ
  ならないとき 200 サイクルであるとする．P1 が再開した後，P2 の実行中に
  キャッシュから追い出された 5\,000 個のデータ項目にアクセスするなら，
  それぞれへの最初のアクセスは $200 - 4 = 196$ サイクルを余分に要し，
  合計で $5\,000\times 196 = 980\,000$ サイクルとなる．間接のコストは
  スイッチの直接のコストのほぼ 1000 倍である．
\end{example}

\begin{note}
  オペレーティングシステムは，キャッシュの状態にかかわらずコンテキストを
  切り替える．例えば，より優先度の高いプロセスが実行可能になったときや，
  実行中のプロセスが割り当てられた時間を使い切ったとき
  （\cref{sec:l02-scheduler}）である．スイッチが残すコールドキャッシュ
  は，CPU を共有するために払う代償の一部である．
\end{note}

\begin{keypoint}
  PCB は，オペレーティングシステムがプロセスを止めて後で再開するために
  必要なすべてを保持する．コンテキストスイッチは，あるプロセスの CPU
  状態をその PCB へ退避し，別のプロセスの状態をその PCB から復元する．
  その隠れたコスト，すなわち残されるコールドキャッシュは，退避と復元の
  コストをはるかに上回りうる．\margin{Li ら \cite{li2007} の測定では，
  データにアクセスしないプロセスの間のスイッチは
  \SI{3.8}{\micro\second} であり，データがキャッシュに収まらなくなると
  合計で \SI{1000}{\micro\second} を超える．}
\end{keypoint}

\section{プロセスのライフサイクル}
\label{sec:l02-states}

CPU から見ると，プロセスは実行中と非実行中を交互に繰り返す．
\source{P2L1，プロセスのライフサイクル，プロセスのライフサイクル: 状態；
OSTEP ch.~4．}割り込まれて切り替えられた実行中のプロセスは休止する．
それは依然として実行の準備ができており，スケジューラが CPU へ
ディスパッチすれば再び実行される．プロセスが休止する理由は他にもあり，
オペレーティングシステムはこれらの状況を
\term[ぷろせすじょうたい]{プロセス状態}[process state]によって区別し，
PCB に記録する．プロセスは次のいずれかの状態にある
（\cref{fig:l02-states}）．

\begin{description}
  \item[新規] 生成されつつあるプロセスは
    \term[しんきじょうたい]{新規状態}[new state]にある．オペレーティング
    システムは\emph{受け入れ制御}を行い，そのプロセスを受け入れられるか
    を判断する．受け入れる場合は PCB といくらかの初期資源を割り当てて
    初期化し，プロセスは\emph{受け入れ}られる．
  \item[実行可能] \term[じっこうかのうじょうたい]{実行可能状態}[ready state]%
    のプロセスは，CPU 以外に実行に必要なものをすべて備えている．
    スケジューラが選びさえすれば実行できるのであり，それまでは順番を
    待つ．
  \item[実行] \term[じっこうじょうたい]{実行状態}[running state]の
    プロセスは CPU 上で実行されている．\emph{スケジューラによる
    ディスパッチ}の後，PCB に格納された地点から実行を続ける．実行中の
    プロセスがこの状態を離れるのは次の3通りである．割り込まれて---例えば CPU
    上の時間を使い切って---コンテキストスイッチの後に実行可能状態へ
    戻る．ディスクからの読み出しやキーボード・タイマのイベント待ちの
    ような長い操作を始めて待ち状態に入る．あるいは終了する．
  \item[待ち] \term[まちじょうたい]{待ち状態}[waiting state]のプロセス
    は，I/O 操作の完了などの何らかの事象が起こるまで先へ進めない．%
    \caution{事象の完了はプロセスを実行可能にするので
    あって，実行中にはしない．CPU は，直ちに行われることもある
    ディスパッチによって与えられる．}事象が起これば，プロセスは再び実行
    可能になる．
  \item[終了] すべての操作を終えたプロセス，あるいは誤りを検出した
    プロセスは \emph{exit} を呼び，成功か誤りかを示す終了コードを返す．
    不正なメモリアクセスなどの障害を引き起こしたプロセスは，
    オペレーティングシステムによって終了させられる．いずれの場合も
    プロセスは\term[しゅうりょうじょうたい]{終了状態}[terminated state]%
    に入る．
\end{description}

\begin{figure}[htbp]
  \centering
  % Process life cycle: states and transitions.
% Usage: % Process life cycle: states and transitions.
% Usage: % Process life cycle: states and transitions.
% Usage: \input{figures/l02-states}   (width ~116 mm)
\begin{tikzpicture}[x=1mm, y=1mm, font=\footnotesize\sffamily,
  st/.style={draw, rounded corners=6pt, fill=ln-box, minimum width=22mm,
             minimum height=8mm, inner sep=2pt},
  >={Stealth[length=1.8mm]},
  lab/.style={font=\scriptsize\sffamily, align=center}]
  \node[st] (new)  at (11,32)  {\iflangja{新規}{new}};
  \node[st] (term) at (105,32) {\iflangja{終了}{terminated}};
  \node[st] (rdy)  at (32,12)  {\iflangja{実行可能}{ready}};
  \node[st, fill=ln-accent!22] (run) at (84,12) {\iflangja{実行}{running}};
  \node[st] (wait) at (58,-12) {\iflangja{待ち}{waiting}};
  \draw[->] (new) -- node[lab, left, pos=0.45] {\iflangja{受け入れ}{admitted}} (rdy);
  \draw[->] (rdy.north east) to[bend left=18]
    node[lab, above] {\iflangja{スケジューラによるディスパッチ}{scheduler dispatch}} (run.north west);
  \draw[->] (run.south west) to[bend left=18]
    node[lab, below] {\iflangja{割り込み}{interrupt}} (rdy.south east);
  \draw[->] (run) -- node[lab, right, pos=0.55] {\iflangja{終了（exit）}{exit}} (term);
  \draw[->] (run) -- node[lab, right, pos=0.6] {\iflangja{I/O・イベント待ち}{I/O or event wait}} (wait);
  \draw[->] (wait) -- node[lab, left, pos=0.4] {\iflangja{I/O・イベント完了}{I/O or event completion}} (rdy);
\end{tikzpicture}
   (width ~116 mm)
\begin{tikzpicture}[x=1mm, y=1mm, font=\footnotesize\sffamily,
  st/.style={draw, rounded corners=6pt, fill=ln-box, minimum width=22mm,
             minimum height=8mm, inner sep=2pt},
  >={Stealth[length=1.8mm]},
  lab/.style={font=\scriptsize\sffamily, align=center}]
  \node[st] (new)  at (11,32)  {\iflangja{新規}{new}};
  \node[st] (term) at (105,32) {\iflangja{終了}{terminated}};
  \node[st] (rdy)  at (32,12)  {\iflangja{実行可能}{ready}};
  \node[st, fill=ln-accent!22] (run) at (84,12) {\iflangja{実行}{running}};
  \node[st] (wait) at (58,-12) {\iflangja{待ち}{waiting}};
  \draw[->] (new) -- node[lab, left, pos=0.45] {\iflangja{受け入れ}{admitted}} (rdy);
  \draw[->] (rdy.north east) to[bend left=18]
    node[lab, above] {\iflangja{スケジューラによるディスパッチ}{scheduler dispatch}} (run.north west);
  \draw[->] (run.south west) to[bend left=18]
    node[lab, below] {\iflangja{割り込み}{interrupt}} (rdy.south east);
  \draw[->] (run) -- node[lab, right, pos=0.55] {\iflangja{終了（exit）}{exit}} (term);
  \draw[->] (run) -- node[lab, right, pos=0.6] {\iflangja{I/O・イベント待ち}{I/O or event wait}} (wait);
  \draw[->] (wait) -- node[lab, left, pos=0.4] {\iflangja{I/O・イベント完了}{I/O or event completion}} (rdy);
\end{tikzpicture}
   (width ~116 mm)
\begin{tikzpicture}[x=1mm, y=1mm, font=\footnotesize\sffamily,
  st/.style={draw, rounded corners=6pt, fill=ln-box, minimum width=22mm,
             minimum height=8mm, inner sep=2pt},
  >={Stealth[length=1.8mm]},
  lab/.style={font=\scriptsize\sffamily, align=center}]
  \node[st] (new)  at (11,32)  {\iflangja{新規}{new}};
  \node[st] (term) at (105,32) {\iflangja{終了}{terminated}};
  \node[st] (rdy)  at (32,12)  {\iflangja{実行可能}{ready}};
  \node[st, fill=ln-accent!22] (run) at (84,12) {\iflangja{実行}{running}};
  \node[st] (wait) at (58,-12) {\iflangja{待ち}{waiting}};
  \draw[->] (new) -- node[lab, left, pos=0.45] {\iflangja{受け入れ}{admitted}} (rdy);
  \draw[->] (rdy.north east) to[bend left=18]
    node[lab, above] {\iflangja{スケジューラによるディスパッチ}{scheduler dispatch}} (run.north west);
  \draw[->] (run.south west) to[bend left=18]
    node[lab, below] {\iflangja{割り込み}{interrupt}} (rdy.south east);
  \draw[->] (run) -- node[lab, right, pos=0.55] {\iflangja{終了（exit）}{exit}} (term);
  \draw[->] (run) -- node[lab, right, pos=0.6] {\iflangja{I/O・イベント待ち}{I/O or event wait}} (wait);
  \draw[->] (wait) -- node[lab, left, pos=0.4] {\iflangja{I/O・イベント完了}{I/O or event completion}} (rdy);
\end{tikzpicture}

  \caption{プロセスの状態と，その間の遷移．}
  \label{fig:l02-states}
\end{figure}

\cref{tab:l02-trace} は，1つのプロセスをそのライフサイクルに沿って
たどったものである．\caution{Linux は実行可能と実行を区別しない．どちらも
\texttt{TASK\_RUNNING} であり，\texttt{ps} はどちらも R と表示する．S と
D は割り込み可能・不可能の2種類の待ちであり，T は停止したプロセスである．}%
プロセスに CPU を与える遷移はディスパッチだけで
あり，他のすべての遷移は CPU を取り上げるか，CPU を持たないままに
することに注意する．

\begin{table}[htbp]
  \centering
  \small
  \caption{実行され，一度割り込まれ，ファイルを読み，終了するプロセスの
    状態の推移．待つことになるキューと，レディキューから取り出す
    ディスパッチについては \cref{sec:l02-scheduler} で述べる．}
  \label{tab:l02-trace}
  \begin{tabular}{@{}llll@{}}
    \toprule
    事象 & 遷移 & 遷移後の状態 & 所在 \\
    \midrule
    プロセスが生成される     & ---                   & 新規     & --- \\
    PCB を初期化し受け入れ   & 受け入れ              & 実行可能 & レディキュー \\
    スケジューラが選ぶ       & ディスパッチ          & 実行     & CPU \\
    CPU 上の時間が満了       & 割り込み              & 実行可能 & レディキュー \\
    スケジューラが選ぶ       & ディスパッチ          & 実行     & CPU \\
    ディスク読み出しを発行   & I/O・イベント待ち     & 待ち     & ディスク I/O キュー \\
    ディスク読み出しが完了   & I/O・イベント完了     & 実行可能 & レディキュー \\
    スケジューラが選ぶ       & ディスパッチ          & 実行     & CPU \\
    \texttt{exit} を呼ぶ     & 終了（exit）          & 終了     & --- \\
    \bottomrule
  \end{tabular}
\end{table}

\section{プロセスの生成}
\label{sec:l02-creation}

プロセスは他のプロセスによって生成される．\source{P2L1，プロセスの
ライフサイクル: 生成；OSTEP ch.~5；Silberschatz ら ch.~3．}システムが
起動してオペレーティングシステムがロードされると，いくつかの初期
プロセスが生成され，その一部は特権的なアクセスを持つ．ユーザがログイン
すると，そのユーザのためにシェルプロセスが生成され，シェルから起動された
プログラム，例えば \texttt{ls} や \texttt{emacs} は，シェルが生成した
新しいプロセスとして実行される．別のプロセスを生成したプロセスは，その
プロセスの\term[おやぷろせす]{親プロセス}[parent process]であり，生成
された側のプロセスは，生成した側の
\term[こぷろせす]{子プロセス}[child process]である．最初のプロセスを
除くすべてのプロセスが親を持つので，システムのプロセスは木をなす
（\cref{fig:l02-creation}\,(a)）．Unix 系システムではその根は
\term[init]{\texttt{init}}[init]，すなわちブート後にカーネルが最初に
開始するユーザプロセスであり，他のすべてのユーザプロセスはそこから派生
する．
（現在の多くの Linux ディストリビューションでは，この最初のプロセスと
して動くプログラムは \texttt{systemd} である．）

\begin{figure}[htbp]
  \centering
  % Left: a process tree.  Right: fork followed by exec.
% Usage: % Left: a process tree.  Right: fork followed by exec.
% Usage: % Left: a process tree.  Right: fork followed by exec.
% Usage: \input{figures/l02-creation}   (width ~116 mm)
\begin{tikzpicture}[x=1mm, y=1mm, font=\footnotesize\sffamily,
  pr/.style={draw, rounded corners=2pt, fill=ln-box, minimum width=13mm,
             minimum height=6mm, inner sep=1pt},
  bx/.style={draw, rounded corners=2pt, fill=white, text width=25mm, align=center,
             minimum height=9mm, inner sep=1.5pt, font=\scriptsize\sffamily},
  >={Stealth[length=1.6mm]},
  lab/.style={font=\scriptsize\sffamily, align=center},
  ttl/.style={font=\footnotesize\sffamily\itshape}]
  % ---- (a) process tree
  \node[pr, fill=ln-accent!22] (init) at (24,36) {\texttt{init}};
  \node[pr] (sh)  at (12,22) {\iflangja{シェル}{shell}};
  \node[pr] (d)   at (36,22) {\iflangja{デーモン}{daemon}};
  \node[pr] (ls)  at (7,8)   {\texttt{ls}};
  \node[pr] (em)  at (22,8)  {\texttt{emacs}};
  \draw[->] (init) -- (sh);
  \draw[->] (init) -- (d);
  \draw[->] (sh) -- (ls);
  \draw[->] (sh) -- (em);
  \node[ttl] at (21,-3) {\iflangja{(a) プロセスの木}{(a) process tree}};
  % ---- (b) fork, then exec in the child
  \node[bx] (par) at (66,34)
    {\iflangja{親：プログラム A\\PC = fork の直後}{parent: program A\\PC = after \texttt{fork}}};
  \node[bx] (ch)  at (66,9)
    {\iflangja{子：A の複製\\PC = fork の直後}{child: copy of A\\PC = after \texttt{fork}}};
  \node[bx] (ch2) at (103,9)
    {\iflangja{子：プログラム B\\PC = B の先頭命令}{child: program B\\PC = first instr.\ of B}};
  \draw[->, thick] (par) -- node[lab, right] {\texttt{fork}} (ch);
  \draw[->, thick] (ch) -- node[lab, above] {\texttt{exec}} (ch2);
  \draw[->, thick] (par.east) -- node[lab, above] {\iflangja{親は実行を続ける}{parent continues}} (117,34);
  \node[ttl] at (84,-3) {\iflangja{(b) 新しいプログラムの起動}{(b) starting a new program}};
\end{tikzpicture}
   (width ~116 mm)
\begin{tikzpicture}[x=1mm, y=1mm, font=\footnotesize\sffamily,
  pr/.style={draw, rounded corners=2pt, fill=ln-box, minimum width=13mm,
             minimum height=6mm, inner sep=1pt},
  bx/.style={draw, rounded corners=2pt, fill=white, text width=25mm, align=center,
             minimum height=9mm, inner sep=1.5pt, font=\scriptsize\sffamily},
  >={Stealth[length=1.6mm]},
  lab/.style={font=\scriptsize\sffamily, align=center},
  ttl/.style={font=\footnotesize\sffamily\itshape}]
  % ---- (a) process tree
  \node[pr, fill=ln-accent!22] (init) at (24,36) {\texttt{init}};
  \node[pr] (sh)  at (12,22) {\iflangja{シェル}{shell}};
  \node[pr] (d)   at (36,22) {\iflangja{デーモン}{daemon}};
  \node[pr] (ls)  at (7,8)   {\texttt{ls}};
  \node[pr] (em)  at (22,8)  {\texttt{emacs}};
  \draw[->] (init) -- (sh);
  \draw[->] (init) -- (d);
  \draw[->] (sh) -- (ls);
  \draw[->] (sh) -- (em);
  \node[ttl] at (21,-3) {\iflangja{(a) プロセスの木}{(a) process tree}};
  % ---- (b) fork, then exec in the child
  \node[bx] (par) at (66,34)
    {\iflangja{親：プログラム A\\PC = fork の直後}{parent: program A\\PC = after \texttt{fork}}};
  \node[bx] (ch)  at (66,9)
    {\iflangja{子：A の複製\\PC = fork の直後}{child: copy of A\\PC = after \texttt{fork}}};
  \node[bx] (ch2) at (103,9)
    {\iflangja{子：プログラム B\\PC = B の先頭命令}{child: program B\\PC = first instr.\ of B}};
  \draw[->, thick] (par) -- node[lab, right] {\texttt{fork}} (ch);
  \draw[->, thick] (ch) -- node[lab, above] {\texttt{exec}} (ch2);
  \draw[->, thick] (par.east) -- node[lab, above] {\iflangja{親は実行を続ける}{parent continues}} (117,34);
  \node[ttl] at (84,-3) {\iflangja{(b) 新しいプログラムの起動}{(b) starting a new program}};
\end{tikzpicture}
   (width ~116 mm)
\begin{tikzpicture}[x=1mm, y=1mm, font=\footnotesize\sffamily,
  pr/.style={draw, rounded corners=2pt, fill=ln-box, minimum width=13mm,
             minimum height=6mm, inner sep=1pt},
  bx/.style={draw, rounded corners=2pt, fill=white, text width=25mm, align=center,
             minimum height=9mm, inner sep=1.5pt, font=\scriptsize\sffamily},
  >={Stealth[length=1.6mm]},
  lab/.style={font=\scriptsize\sffamily, align=center},
  ttl/.style={font=\footnotesize\sffamily\itshape}]
  % ---- (a) process tree
  \node[pr, fill=ln-accent!22] (init) at (24,36) {\texttt{init}};
  \node[pr] (sh)  at (12,22) {\iflangja{シェル}{shell}};
  \node[pr] (d)   at (36,22) {\iflangja{デーモン}{daemon}};
  \node[pr] (ls)  at (7,8)   {\texttt{ls}};
  \node[pr] (em)  at (22,8)  {\texttt{emacs}};
  \draw[->] (init) -- (sh);
  \draw[->] (init) -- (d);
  \draw[->] (sh) -- (ls);
  \draw[->] (sh) -- (em);
  \node[ttl] at (21,-3) {\iflangja{(a) プロセスの木}{(a) process tree}};
  % ---- (b) fork, then exec in the child
  \node[bx] (par) at (66,34)
    {\iflangja{親：プログラム A\\PC = fork の直後}{parent: program A\\PC = after \texttt{fork}}};
  \node[bx] (ch)  at (66,9)
    {\iflangja{子：A の複製\\PC = fork の直後}{child: copy of A\\PC = after \texttt{fork}}};
  \node[bx] (ch2) at (103,9)
    {\iflangja{子：プログラム B\\PC = B の先頭命令}{child: program B\\PC = first instr.\ of B}};
  \draw[->, thick] (par) -- node[lab, right] {\texttt{fork}} (ch);
  \draw[->, thick] (ch) -- node[lab, above] {\texttt{exec}} (ch2);
  \draw[->, thick] (par.east) -- node[lab, above] {\iflangja{親は実行を続ける}{parent continues}} (117,34);
  \node[ttl] at (84,-3) {\iflangja{(b) 新しいプログラムの起動}{(b) starting a new program}};
\end{tikzpicture}

  \caption{(a) プロセスは初期プロセスを根とする木をなす．(b)
    \texttt{fork} で生成された子は親の複製であり，\texttt{exec} が
    そのプログラムを新しいものに置き換える．}
  \label{fig:l02-creation}
\end{figure}

ほとんどのオペレーティングシステムは，\caution{\texttt{fork} がコピーするのはメモリでは
なくページテーブルである．ページが複製されるのは，どちらかが書き込む
ときだけである（コピーオンライト，第8章）．}%
プロセスを生成する2つの機構を
提供する．Unix ではそれらを \texttt{fork} と \texttt{exec} と呼ぶ．
\begin{description}
  \item[\texttt{fork}] オペレーティングシステムは，
    \term*{\texttt{fork}}[fork]\index{fork@\texttt{fork}} システムコールによって子のために
    新しい PCB を作り，親の PCB をそこへコピーし，親のアドレス空間の
    複製を子に与える．子は親の複製であり，同じプログラムと同じプログラム
    カウンタを持つので，呼び出しの後，親も子も \texttt{fork} の次の命令
    から実行を続ける．
  \item[\texttt{exec}] プロセスは，\term[exec]{\texttt{exec}}[exec]
    システムコールによって自分のイメージを新しいプログラムで置き換える．
    オペレーティングシステムはプロセスのアドレス空間を新しいプログラム
    で置き換え，それを記述するように PCB を更新する．プログラム
    カウンタは新しいプログラムの最初の命令を指すようになり，プロセス
    識別子とオープンしたファイルは保たれる．
\end{description}

\emph{異なる}プログラムを実行するプロセスを生成するには，親が
\texttt{fork} を呼び，続いて子が \texttt{exec} を呼ぶ
（\cref{fig:l02-creation}\,(b)）．親と子は同じ \texttt{fork} から同じ
場所へ戻るので，その戻り値によって互いを区別する．\texttt{fork} は，子
では 0 を，親では子のプロセス識別子を返し，子を生成できなかった場合は親
で $-1$ を返す．

Android は，アプリケーションを素早く起動するために \texttt{fork} を
用いる．そのアプリケーションプロセスは，ブート時に開始され，
アプリケーションの起動だけを目的とするデーモンである
\term{Zygote}[Zygote] から派生する．Zygote はランタイムとよく使われる
ライブラリを一度だけロードしておく．\sidenote{\texttt{fork} は親の
ページをコピーオンライトで共有するので，ランタイムとあらかじめロード
されたクラスは各アプリケーションへコピーされない．どのアプリケーション
プロセスも，それらがすでに対応付けられた状態で始まり，誰も書き込まない
限り物理ページは1組のままである．Zygote はこうして起動時間とメモリを
同時に節約する．}アプリケーションを起動するときには
Zygote が fork し，Zygote がロードしたものをすでにすべて含む子が，
そのアプリケーションのプロセスとなる．

\Needspace{18\baselineskip}
\begin{example}
  \texttt{fork} で子プロセスを生成し，その中で \texttt{exec} により
  \texttt{ls -l} を実行し，子が終わるまで待つプログラムを示す．
  \begin{lstlisting}[language=C, numbers=none]
#include <sys/types.h>
#include <sys/wait.h>
#include <unistd.h>

int main(void) {
    pid_t pid = fork();
    if (pid == 0) {              /* 子 */
        execlp("ls", "ls", "-l", (char *)NULL);
        _exit(127);              /* exec 失敗時のみ */
    } else if (pid > 0) {        /* 親 */
        waitpid(pid, NULL, 0);   /* 子の終了を待つ */
    }
    return 0;
}
\end{lstlisting}
  \texttt{execlp} が成功すれば，子は \texttt{ls} のコードを実行する．
  この呼び出しが戻ることはない．\margin{exit したプロセスは，親が終了
  コードを読み取る---ここでは \texttt{waitpid} で---まで PCB を保持し
  続ける．それまではゾンビ（状態 Z）である．親が先に終了した子は
  \texttt{init} に引き取られ回収される．}
\end{example}

\section{CPU スケジューラ}
\label{sec:l02-scheduler}

CPU がプロセスの実行を開始するには，そのプロセスが実行可能でなければ
ならず，実行可能なプロセスはいつでも複数ありうる．\source{P2L1，CPU
スケジューラの役割；OSTEP ch.~6--7．}オペレーティングシステムは，その
うちのどれを次に実行するか，そしてどれだけの間実行させるかを決めなければ
ならない．

\begin{definition}[CPUスケジューラ]
  \term[CPUすけじゅーら]{CPUスケジューラ}[CPU scheduler]とは，プロセスが
  CPU をどのように使うかを管理するオペレーティングシステムの構成要素で
  ある．実行可能なプロセスのどれを CPU へディスパッチするか，そして
  どれだけの間実行させてよいかを決める．
\end{definition}

実行可能なプロセスは\term[れでぃきゅー]{レディキュー}[ready queue]で
CPU を待つ．そのうちの1つに CPU を与えるために，オペレーティング
システムは3つの操作を行う．
\begin{description}
  \item[プリエンプト] 実行中のプロセスに割り込み，そのコンテキストを
    PCB へ退避する．これを
    \term[ぷりえんぷしょん]{プリエンプション}[preemption]と呼ぶ．
  \item[スケジュール] スケジューラがアルゴリズムを実行し，レディ
    キューから次のプロセスを選ぶ．
  \item[ディスパッチ] オペレーティングシステムは選ばれたプロセスの
    コンテキストへ切り替え，CPU 上で実行させる．これがプロセスの
    \term[でぃすぱっち]{ディスパッチ}[dispatch]である．
\end{description}

CPU はプロセスが必要とする資源であるから，その時間はできるだけ，
オペレーティングシステムの実行ではなくプロセスの実行に費やされるべきで
ある．そのためスケジューラには，効率的なアルゴリズムと，その効率的な
実装，そして待機中のプロセスとそのスケジューリング情報を表す効率的な
データ構造---とりわけレディキュー---が必要である．

\subsection{スケジューラを動かす頻度}

スケジューラを動かす回数が多いほど，プロセスから奪う CPU 時間は増える．
\source{P2L1，プロセスの長さ；OSTEP ch.~7．}各プロセスが時間 $T_p$ だけ
実行された後，$t_{\mathrm{sched}}$ を要するスケジューラが次のプロセスを
選ぶとする（\cref{fig:l02-timeslice}）．これを $n$ 回繰り返すと，
有用な処理に費やされる CPU 時間の割合は次のようになる．
\begin{equation}
  \text{有用な CPU 処理}
  \;=\; \frac{\text{総処理時間}}{\text{総時間}}
  \;=\; \frac{n\,T_p}{n\,T_p + n\,t_{\mathrm{sched}}}
  \;=\; \frac{T_p}{T_p + t_{\mathrm{sched}}} .
  \label{eq:l02-useful}
\end{equation}

\begin{figure}[htbp]
  \centering
  % Alternation of process timeslices and scheduler runs.
% Usage: % Alternation of process timeslices and scheduler runs.
% Usage: % Alternation of process timeslices and scheduler runs.
% Usage: \input{figures/l02-timeslice}   (width ~117 mm incl. the CPU label)
\begin{tikzpicture}[x=1mm, y=1mm, font=\footnotesize\sffamily,
  dim/.style={{Bar[width=2mm]}-{Bar[width=2mm]}, thin},
  pr/.style={draw, fill=ln-accent!22},
  sc/.style={draw, fill=ln-box}]
  \node[anchor=east] at (-1,3) {CPU};
  \draw[pr] (0,0)  rectangle (30,6)  node[midway] {P1};
  \draw[sc] (30,0) rectangle (38,6)  node[midway, font=\scriptsize\sffamily] {\iflangja{S}{S}};
  \draw[pr] (38,0) rectangle (68,6)  node[midway] {P2};
  \draw[sc] (68,0) rectangle (76,6)  node[midway, font=\scriptsize\sffamily] {\iflangja{S}{S}};
  \draw[pr] (76,0) rectangle (106,6) node[midway] {P3};
  \draw[sc] (106,0) rectangle (109,6);
  \foreach \a/\b in {0/30, 38/68, 76/106} {
    \draw[dim] (\a,-3) -- (\b,-3);
    \node[below] at ({(\a+\b)/2},-3.5) {$T_p$};
  }
  \foreach \a/\b in {30/38, 68/76} {
    \draw[dim] (\a,9) -- (\b,9);
    \node[above] at ({(\a+\b)/2},9.5) {$t_{\mathrm{sched}}$};
  }
\end{tikzpicture}
   (width ~117 mm incl. the CPU label)
\begin{tikzpicture}[x=1mm, y=1mm, font=\footnotesize\sffamily,
  dim/.style={{Bar[width=2mm]}-{Bar[width=2mm]}, thin},
  pr/.style={draw, fill=ln-accent!22},
  sc/.style={draw, fill=ln-box}]
  \node[anchor=east] at (-1,3) {CPU};
  \draw[pr] (0,0)  rectangle (30,6)  node[midway] {P1};
  \draw[sc] (30,0) rectangle (38,6)  node[midway, font=\scriptsize\sffamily] {\iflangja{S}{S}};
  \draw[pr] (38,0) rectangle (68,6)  node[midway] {P2};
  \draw[sc] (68,0) rectangle (76,6)  node[midway, font=\scriptsize\sffamily] {\iflangja{S}{S}};
  \draw[pr] (76,0) rectangle (106,6) node[midway] {P3};
  \draw[sc] (106,0) rectangle (109,6);
  \foreach \a/\b in {0/30, 38/68, 76/106} {
    \draw[dim] (\a,-3) -- (\b,-3);
    \node[below] at ({(\a+\b)/2},-3.5) {$T_p$};
  }
  \foreach \a/\b in {30/38, 68/76} {
    \draw[dim] (\a,9) -- (\b,9);
    \node[above] at ({(\a+\b)/2},9.5) {$t_{\mathrm{sched}}$};
  }
\end{tikzpicture}
   (width ~117 mm incl. the CPU label)
\begin{tikzpicture}[x=1mm, y=1mm, font=\footnotesize\sffamily,
  dim/.style={{Bar[width=2mm]}-{Bar[width=2mm]}, thin},
  pr/.style={draw, fill=ln-accent!22},
  sc/.style={draw, fill=ln-box}]
  \node[anchor=east] at (-1,3) {CPU};
  \draw[pr] (0,0)  rectangle (30,6)  node[midway] {P1};
  \draw[sc] (30,0) rectangle (38,6)  node[midway, font=\scriptsize\sffamily] {\iflangja{S}{S}};
  \draw[pr] (38,0) rectangle (68,6)  node[midway] {P2};
  \draw[sc] (68,0) rectangle (76,6)  node[midway, font=\scriptsize\sffamily] {\iflangja{S}{S}};
  \draw[pr] (76,0) rectangle (106,6) node[midway] {P3};
  \draw[sc] (106,0) rectangle (109,6);
  \foreach \a/\b in {0/30, 38/68, 76/106} {
    \draw[dim] (\a,-3) -- (\b,-3);
    \node[below] at ({(\a+\b)/2},-3.5) {$T_p$};
  }
  \foreach \a/\b in {30/38, 68/76} {
    \draw[dim] (\a,9) -- (\b,9);
    \node[above] at ({(\a+\b)/2},9.5) {$t_{\mathrm{sched}}$};
  }
\end{tikzpicture}

  \caption{各プロセスは $T_p$ ずつ実行され，その間にスケジューラ（S）が
    $t_{\mathrm{sched}}$ だけ実行される．}
  \label{fig:l02-timeslice}
\end{figure}

プロセスが CPU 上で割り当てられる時間 $T_p$ が，そのプロセスの
\term[たいむすらいす]{タイムスライス}[timeslice]である．
$t_{\mathrm{sched}}$ に対してタイムスライスが長いほど，有用な割合は
大きくなる．タイムスライスがスケジューラ1回分の長さしかなければ CPU
時間の半分が失われ，その 10 倍の長さのタイムスライスなら有用な割合は
$10/11 \approx \SI{91}{\percent}$ となる．

\Needspace{8\baselineskip}
\begin{example}
  $t_{\mathrm{sched}} = \SI{0.5}{\milli\second}$ とする．タイムスライス
  が \SI{2}{\milli\second} なら，\cref{eq:l02-useful} より有用な処理は
  $2/2.5 = \SI{80}{\percent}$ であり，タイムスライスが
  \SI{20}{\milli\second} なら $20/20.5 \approx \SI{97.6}{\percent}$ と
  なる．
\end{example}

しかし，長いタイムスライスにも代償がある．1つのプロセスが CPU を握って
いる間，他の実行可能なプロセスは順番を待つ．したがってスケジューラを
選ぶには設計上の判断が要る．どのようなタイムスライスの値が適切か，
そしてどの指標によって次に実行するプロセスを決めるかである．第7章では
これらの問いを扱う．\margin{実際の値もこの程度である．Linux の CFS は
\SI{6}{\milli\second} の周期を実行可能なタスクで分け合わせ，1タスク
あたり \SI{0.75}{\milli\second} を下限とした．Linux~6.6 の EEVDF は
\SI{0.7}{\milli\second} の基本スライスを用いる．}

\subsection{プロセスと I/O}

実行中のプロセスは，例えばディスクから読み出すために I/O 要求を発行する
ことがある．\source{P2L1，I/O はどうなるか；Silberschatz ら ch.~3．}%
オペレーティングシステムはその要求をデバイスへ届け，プロセスをその
デバイスに対応する \emph{I/O キュー}に置く．プロセスは，デバイスが操作を
完了して要求に応答するまでそこで待つ．その後プロセスは再び実行可能に
なる．システムの負荷に応じてレディキューに置かれるか，他に待っている
プロセスがなければ直ちに CPU へディスパッチされることもある．%
\margin{この待ちは，この章の他のすべてと比べて長い．\SI{4}{\kibi\byte}
の読み出しには，NVMe SSD で数十マイクロ秒，回転ディスクで
\SI{10}{\milli\second} 程度を要する．}

\cref{fig:l02-queues} は，プロセスがレディキューに入る，あるいは再び
入る経路を示している．
\begin{itemize}
  \item I/O 要求がデバイスによって処理された．
  \item タイムスライスが満了した．
  \item \texttt{fork} によって新しいプロセスが生成された．子がレディ
    キューに置かれる．
  \item 割り込みを待っていて，その割り込みが発生した．%
    \margin{割り込みの処理自体はコンテキストスイッチではない．ハンドラ
    は，割り込まれたプロセスのコンテキストのまま，カーネル内で実行される
    （第5章）．}
\end{itemize}

\begin{figure}[htbp]
  \centering
  % Queueing diagram: how processes move between the ready queue, the CPU
% and I/O queues.
% Usage: % Queueing diagram: how processes move between the ready queue, the CPU
% and I/O queues.
% Usage: % Queueing diagram: how processes move between the ready queue, the CPU
% and I/O queues.
% Usage: \input{figures/l02-queues}   (width ~108 mm)
\begin{tikzpicture}[x=1mm, y=1mm, font=\footnotesize\sffamily,
  >={Stealth[length=1.6mm]},
  res/.style={draw, circle, fill=ln-accent!22, minimum size=10mm, inner sep=0pt},
  bx/.style={draw, rounded corners=2pt, fill=white, minimum height=6mm,
             minimum width=22mm, inner sep=1pt, font=\scriptsize\sffamily},
  lab/.style={font=\scriptsize\sffamily}]
  % queue drawing: rectangle with slots, open at the left
  \newcommand{\lnqueue}[3]{% x, y, label
    \draw[fill=ln-box] (#1,#2-3) rectangle ++(28,6);
    \foreach \i in {1,...,4} \draw (#1+28-4*\i,#2-3) -- ++(0,6);
    \node[lab, anchor=south] at (#1+14,#2+3) {#3};}
  \lnqueue{16}{40}{\iflangja{レディキュー}{ready queue}}
  \node[res] (cpu) at (64,40) {CPU};
  \draw[->] (44,40) -- (cpu.west);
  % right-hand bus and left-hand return line
  \draw (cpu.east) -- (102,40) -- (102,-8);
  \draw[->] (2,-8) -- (2,40) -- (16,40);
  % I/O
  \draw[->] (102,26) -- node[lab, above] {\iflangja{I/O 要求}{I/O request}} (80,26);
  \lnqueue{52}{26}{\iflangja{I/O キュー}{I/O queue}}
  \node[res, minimum size=8mm] (io) at (36,26) {I/O};
  \draw[->] (52,26) -- (io.east);
  \draw (io.west) -- (2,26);
  % timeslice expired
  \draw (102,14) -- node[lab, above, pos=0.5] {\iflangja{タイムスライス満了}{timeslice expired}} (2,14);
  % fork
  \draw[->] (102,3) -- node[lab, above] {\iflangja{子を fork}{fork a child}} (74,3);
  \node[bx] (ch) at (63,3) {\iflangja{子プロセスが作られる}{child is created}};
  \draw (ch.west) -- (2,3);
  % interrupt
  \draw[->] (102,-8) -- node[lab, above] {\iflangja{割り込みを待つ}{wait for an interrupt}} (74,-8);
  \node[bx] (in) at (63,-8) {\iflangja{割り込みが発生}{interrupt occurs}};
  \draw (in.west) -- (2,-8);
\end{tikzpicture}
   (width ~108 mm)
\begin{tikzpicture}[x=1mm, y=1mm, font=\footnotesize\sffamily,
  >={Stealth[length=1.6mm]},
  res/.style={draw, circle, fill=ln-accent!22, minimum size=10mm, inner sep=0pt},
  bx/.style={draw, rounded corners=2pt, fill=white, minimum height=6mm,
             minimum width=22mm, inner sep=1pt, font=\scriptsize\sffamily},
  lab/.style={font=\scriptsize\sffamily}]
  % queue drawing: rectangle with slots, open at the left
  \newcommand{\lnqueue}[3]{% x, y, label
    \draw[fill=ln-box] (#1,#2-3) rectangle ++(28,6);
    \foreach \i in {1,...,4} \draw (#1+28-4*\i,#2-3) -- ++(0,6);
    \node[lab, anchor=south] at (#1+14,#2+3) {#3};}
  \lnqueue{16}{40}{\iflangja{レディキュー}{ready queue}}
  \node[res] (cpu) at (64,40) {CPU};
  \draw[->] (44,40) -- (cpu.west);
  % right-hand bus and left-hand return line
  \draw (cpu.east) -- (102,40) -- (102,-8);
  \draw[->] (2,-8) -- (2,40) -- (16,40);
  % I/O
  \draw[->] (102,26) -- node[lab, above] {\iflangja{I/O 要求}{I/O request}} (80,26);
  \lnqueue{52}{26}{\iflangja{I/O キュー}{I/O queue}}
  \node[res, minimum size=8mm] (io) at (36,26) {I/O};
  \draw[->] (52,26) -- (io.east);
  \draw (io.west) -- (2,26);
  % timeslice expired
  \draw (102,14) -- node[lab, above, pos=0.5] {\iflangja{タイムスライス満了}{timeslice expired}} (2,14);
  % fork
  \draw[->] (102,3) -- node[lab, above] {\iflangja{子を fork}{fork a child}} (74,3);
  \node[bx] (ch) at (63,3) {\iflangja{子プロセスが作られる}{child is created}};
  \draw (ch.west) -- (2,3);
  % interrupt
  \draw[->] (102,-8) -- node[lab, above] {\iflangja{割り込みを待つ}{wait for an interrupt}} (74,-8);
  \node[bx] (in) at (63,-8) {\iflangja{割り込みが発生}{interrupt occurs}};
  \draw (in.west) -- (2,-8);
\end{tikzpicture}
   (width ~108 mm)
\begin{tikzpicture}[x=1mm, y=1mm, font=\footnotesize\sffamily,
  >={Stealth[length=1.6mm]},
  res/.style={draw, circle, fill=ln-accent!22, minimum size=10mm, inner sep=0pt},
  bx/.style={draw, rounded corners=2pt, fill=white, minimum height=6mm,
             minimum width=22mm, inner sep=1pt, font=\scriptsize\sffamily},
  lab/.style={font=\scriptsize\sffamily}]
  % queue drawing: rectangle with slots, open at the left
  \newcommand{\lnqueue}[3]{% x, y, label
    \draw[fill=ln-box] (#1,#2-3) rectangle ++(28,6);
    \foreach \i in {1,...,4} \draw (#1+28-4*\i,#2-3) -- ++(0,6);
    \node[lab, anchor=south] at (#1+14,#2+3) {#3};}
  \lnqueue{16}{40}{\iflangja{レディキュー}{ready queue}}
  \node[res] (cpu) at (64,40) {CPU};
  \draw[->] (44,40) -- (cpu.west);
  % right-hand bus and left-hand return line
  \draw (cpu.east) -- (102,40) -- (102,-8);
  \draw[->] (2,-8) -- (2,40) -- (16,40);
  % I/O
  \draw[->] (102,26) -- node[lab, above] {\iflangja{I/O 要求}{I/O request}} (80,26);
  \lnqueue{52}{26}{\iflangja{I/O キュー}{I/O queue}}
  \node[res, minimum size=8mm] (io) at (36,26) {I/O};
  \draw[->] (52,26) -- (io.east);
  \draw (io.west) -- (2,26);
  % timeslice expired
  \draw (102,14) -- node[lab, above, pos=0.5] {\iflangja{タイムスライス満了}{timeslice expired}} (2,14);
  % fork
  \draw[->] (102,3) -- node[lab, above] {\iflangja{子を fork}{fork a child}} (74,3);
  \node[bx] (ch) at (63,3) {\iflangja{子プロセスが作られる}{child is created}};
  \draw (ch.west) -- (2,3);
  % interrupt
  \draw[->] (102,-8) -- node[lab, above] {\iflangja{割り込みを待つ}{wait for an interrupt}} (74,-8);
  \node[bx] (in) at (63,-8) {\iflangja{割り込みが発生}{interrupt occurs}};
  \draw (in.west) -- (2,-8);
\end{tikzpicture}

  \caption{プロセスがレディキューに入る経路．どの経路もそこで終わる．
    \texttt{fork} の経路でキューに加わるのは，呼び出しを発行した
    プロセスではなく新しい子である．}
  \label{fig:l02-queues}
\end{figure}

I/O キューにあるプロセスは，デバイスが要求を完了するまでスケジューラの
制御の外にある．スケジューラが選ぶのはレディキューにあるプロセスだけで
あり，いつ，どのプロセスへコンテキストを切り替えるかを決める．

\Needspace{7\baselineskip}
\begin{keypoint}
  スケジューラは実行中のプロセスをプリエンプトし，レディキューから次の
  プロセスを選んでディスパッチする．スケジューラは同じ CPU 上で実行
  されるので，それ自身の時間はオーバーヘッドであり
  （\cref{eq:l02-useful}），$t_{\mathrm{sched}}$ に対してタイムスライスが
  長くなるほど小さくなる．
\end{keypoint}

\section{プロセス間通信}
\label{sec:l02-ipc}

多くのアプリケーションは，互いにやり取りしなければならない複数の
プロセスとして構成される．例えばウェブアプリケーションは，
クライアントの要求を処理するウェブサーバプロセスと，データを保存する
データベースプロセスからなりうる．\source{P2L1，プロセスは相互作用
できるか；Silberschatz ら ch.~3．}しかしオペレーティングシステムは
プロセスを互いに隔離する（第1章）．各プロセスは自分のアドレス空間を
持ち，どのプロセスも他のプロセスのメモリにアクセスできない．

\begin{definition}[プロセス間通信]
  プロセスが相互作用するためにオペレーティングシステムが提供する機構を
  \term[ぷろせすかんつうしん]{プロセス間通信}[inter-process communication]%
  \index{IPC|see{プロセス間通信}}（IPC）と呼ぶ．IPC の機構は，アドレス
  空間の間でデータや情報を転送し，プロセスの保護と隔離を保ち，
  アプリケーションに柔軟性と良好な性能を与える．
\end{definition}

IPC には基本となる2つの種類がある（\cref{fig:l02-ipc}）．
\begin{description}
  \item[メッセージパッシング]
    \term[めっせーじぱっしんぐ]{メッセージパッシング}[message passing]%
    では，オペレーティングシステムが通信チャネル，例えばカーネルメモリ
    内の共有バッファを確立する．プロセスはシステムコールを用いて，
    チャネルに書き込むことでメッセージを送り，チャネルから読むことで
    メッセージを受け取る．オペレーティングシステムがチャネルを管理し，
    \texttt{send} や \texttt{recv} のような，明確に定義された同一の
    インタフェースをすべてのアプリケーションに提供する．その代償は
    オーバーヘッドである．あらゆる情報が，送信側のアドレス空間から
    カーネル内のチャネルへ，さらにそこから受信側のアドレス空間へと
    コピーされ，その両側でカーネル境界の横断が生じる．
  \item[共有メモリ] \term[きょうゆうめもり]{共有メモリ}[shared memory]%
    では，オペレーティングシステムが物理メモリの領域を確保し，それを
    両方のプロセスのアドレス空間へ対応付ける．プロセスはこのメモリを，
    自分のメモリと同じように直接読み書きし，オペレーティングシステムは
    介在しない．その代償は，オペレーティングシステムがそのメモリの
    使い方を定めなくなることである．明確に定義された固定の
    インタフェースはなく，プロセスはデータの形式について合意し，
    アクセスを自分たちで調整しなければならない．開発者は
    アプリケーションごとにこのコードを書き直すことになりかねず，誤りも
    生じやすい．
\end{description}

\begin{figure}[htbp]
  \centering
  % Message-passing IPC versus shared-memory IPC.
% Usage: % Message-passing IPC versus shared-memory IPC.
% Usage: % Message-passing IPC versus shared-memory IPC.
% Usage: \input{figures/l02-ipc}   (width ~116 mm)
\begin{tikzpicture}[x=1mm, y=1mm, font=\footnotesize\sffamily,
  >={Stealth[length=1.6mm]},
  lab/.style={font=\scriptsize\sffamily, align=center},
  pname/.style={font=\footnotesize\sffamily\bfseries, anchor=north},
  ttl/.style={font=\footnotesize\sffamily\itshape}]
  % ---- (a) message passing
  \draw[rounded corners=2pt, fill=white] (0,23) rectangle (20,38);
  \node[pname] at (10,38) {P1};
  \draw[rounded corners=2pt, fill=white] (34,23) rectangle (54,38);
  \node[pname] at (44,38) {P2};
  \draw[fill=ln-box] (0,0) rectangle (54,13);
  \node[lab, anchor=south west, text=ln-muted] at (0.5,0.3) {\iflangja{カーネル}{kernel}};
  \draw[fill=ln-accent!22] (16,4) rectangle (38,10) node[midway, lab] {\iflangja{チャネル}{channel}};
  \draw[->, thick] (14,23) -- (20,10);
  \draw[->, thick] (34,10) -- (40,23);
  \node[lab, anchor=east] at (13.5,17.5) {\texttt{send}\\\iflangja{（コピー）}{(copy)}};
  \node[lab, anchor=west] at (40.5,17.5) {\texttt{recv}\\\iflangja{（コピー）}{(copy)}};
  \node[ttl] at (27,-5) {\iflangja{(a) メッセージパッシング}{(a) message passing}};
  % ---- (b) shared memory
  \draw[rounded corners=2pt, fill=white] (62,23) rectangle (82,38);
  \node[pname] at (72,38) {P1};
  \draw[rounded corners=2pt, fill=white] (96,23) rectangle (116,38);
  \node[pname] at (106,38) {P2};
  \draw[fill=ln-accent!22] (66,25) rectangle (78,30);
  \draw[fill=ln-accent!22] (100,28) rectangle (112,33);
  \draw[fill=ln-box] (62,0) rectangle (116,13);
  \node[lab, anchor=south west, text=ln-muted] at (62.5,0.3) {\iflangja{物理メモリ}{physical memory}};
  \draw[fill=ln-accent!22] (80,4) rectangle (98,10) node[midway, lab] {\iflangja{共有領域}{shared}};
  \draw[thick, densely dashed] (72,25) -- (85,10);
  \draw[thick, densely dashed] (106,28) -- (93,10);
  \node[lab, text=ln-muted] at (89,19) {\iflangja{対応付け}{mapped}};
  \node[ttl] at (89,-5) {\iflangja{(b) 共有メモリ}{(b) shared memory}};
\end{tikzpicture}
   (width ~116 mm)
\begin{tikzpicture}[x=1mm, y=1mm, font=\footnotesize\sffamily,
  >={Stealth[length=1.6mm]},
  lab/.style={font=\scriptsize\sffamily, align=center},
  pname/.style={font=\footnotesize\sffamily\bfseries, anchor=north},
  ttl/.style={font=\footnotesize\sffamily\itshape}]
  % ---- (a) message passing
  \draw[rounded corners=2pt, fill=white] (0,23) rectangle (20,38);
  \node[pname] at (10,38) {P1};
  \draw[rounded corners=2pt, fill=white] (34,23) rectangle (54,38);
  \node[pname] at (44,38) {P2};
  \draw[fill=ln-box] (0,0) rectangle (54,13);
  \node[lab, anchor=south west, text=ln-muted] at (0.5,0.3) {\iflangja{カーネル}{kernel}};
  \draw[fill=ln-accent!22] (16,4) rectangle (38,10) node[midway, lab] {\iflangja{チャネル}{channel}};
  \draw[->, thick] (14,23) -- (20,10);
  \draw[->, thick] (34,10) -- (40,23);
  \node[lab, anchor=east] at (13.5,17.5) {\texttt{send}\\\iflangja{（コピー）}{(copy)}};
  \node[lab, anchor=west] at (40.5,17.5) {\texttt{recv}\\\iflangja{（コピー）}{(copy)}};
  \node[ttl] at (27,-5) {\iflangja{(a) メッセージパッシング}{(a) message passing}};
  % ---- (b) shared memory
  \draw[rounded corners=2pt, fill=white] (62,23) rectangle (82,38);
  \node[pname] at (72,38) {P1};
  \draw[rounded corners=2pt, fill=white] (96,23) rectangle (116,38);
  \node[pname] at (106,38) {P2};
  \draw[fill=ln-accent!22] (66,25) rectangle (78,30);
  \draw[fill=ln-accent!22] (100,28) rectangle (112,33);
  \draw[fill=ln-box] (62,0) rectangle (116,13);
  \node[lab, anchor=south west, text=ln-muted] at (62.5,0.3) {\iflangja{物理メモリ}{physical memory}};
  \draw[fill=ln-accent!22] (80,4) rectangle (98,10) node[midway, lab] {\iflangja{共有領域}{shared}};
  \draw[thick, densely dashed] (72,25) -- (85,10);
  \draw[thick, densely dashed] (106,28) -- (93,10);
  \node[lab, text=ln-muted] at (89,19) {\iflangja{対応付け}{mapped}};
  \node[ttl] at (89,-5) {\iflangja{(b) 共有メモリ}{(b) shared memory}};
\end{tikzpicture}
   (width ~116 mm)
\begin{tikzpicture}[x=1mm, y=1mm, font=\footnotesize\sffamily,
  >={Stealth[length=1.6mm]},
  lab/.style={font=\scriptsize\sffamily, align=center},
  pname/.style={font=\footnotesize\sffamily\bfseries, anchor=north},
  ttl/.style={font=\footnotesize\sffamily\itshape}]
  % ---- (a) message passing
  \draw[rounded corners=2pt, fill=white] (0,23) rectangle (20,38);
  \node[pname] at (10,38) {P1};
  \draw[rounded corners=2pt, fill=white] (34,23) rectangle (54,38);
  \node[pname] at (44,38) {P2};
  \draw[fill=ln-box] (0,0) rectangle (54,13);
  \node[lab, anchor=south west, text=ln-muted] at (0.5,0.3) {\iflangja{カーネル}{kernel}};
  \draw[fill=ln-accent!22] (16,4) rectangle (38,10) node[midway, lab] {\iflangja{チャネル}{channel}};
  \draw[->, thick] (14,23) -- (20,10);
  \draw[->, thick] (34,10) -- (40,23);
  \node[lab, anchor=east] at (13.5,17.5) {\texttt{send}\\\iflangja{（コピー）}{(copy)}};
  \node[lab, anchor=west] at (40.5,17.5) {\texttt{recv}\\\iflangja{（コピー）}{(copy)}};
  \node[ttl] at (27,-5) {\iflangja{(a) メッセージパッシング}{(a) message passing}};
  % ---- (b) shared memory
  \draw[rounded corners=2pt, fill=white] (62,23) rectangle (82,38);
  \node[pname] at (72,38) {P1};
  \draw[rounded corners=2pt, fill=white] (96,23) rectangle (116,38);
  \node[pname] at (106,38) {P2};
  \draw[fill=ln-accent!22] (66,25) rectangle (78,30);
  \draw[fill=ln-accent!22] (100,28) rectangle (112,33);
  \draw[fill=ln-box] (62,0) rectangle (116,13);
  \node[lab, anchor=south west, text=ln-muted] at (62.5,0.3) {\iflangja{物理メモリ}{physical memory}};
  \draw[fill=ln-accent!22] (80,4) rectangle (98,10) node[midway, lab] {\iflangja{共有領域}{shared}};
  \draw[thick, densely dashed] (72,25) -- (85,10);
  \draw[thick, densely dashed] (106,28) -- (93,10);
  \node[lab, text=ln-muted] at (89,19) {\iflangja{対応付け}{mapped}};
  \node[ttl] at (89,-5) {\iflangja{(b) 共有メモリ}{(b) shared memory}};
\end{tikzpicture}

  \caption{(a) メッセージパッシングは，すべてのメッセージをカーネル内の
    チャネルへコピーし，そこから再びコピーする．(b) 共有メモリでは，
    同じ物理ページが両方のアドレス空間へ対応付けられる．}
  \label{fig:l02-ipc}
\end{figure}

共有メモリは，通信のコストを個々のやり取りから設定時へ移す．領域が一度
対応付けられれば，1回のやり取りにコピーもカーネル境界の横断も伴わないが，
対応付けの作成自体が高価なので，2つの機構のどちらが安上がりかは，
プロセスが何回やり取りするかで決まる．\margin{Linux での測定では，
パイプを通る往復は \SI{10}{\micro\second} 程度，Unix ドメインソケット
ではそれよりやや大きい（lmbench の \texttt{lat\_pipe}，
\texttt{lat\_unix}）．最小のシステムコール1回だけでも
\SI{100}{\nano\second} 程度を要する．}

\Needspace{7\baselineskip}
\begin{example}
  カーネルのチャネルを通じてメッセージを1つ送るコストが
  \SI{6}{\micro\second}（2回のシステムコールと2回のコピー）であり，
  同じデータを共有メモリでやり取りするコストが，一度限りの
  \SI{300}{\micro\second} で対応付けを確立した後は
  \SI{1}{\micro\second} であるとする．$n$ 個のメッセージに対し，共有
  メモリのコストは $300 + n$ マイクロ秒，メッセージパッシングのコストは
  $6n$ マイクロ秒である．共有メモリが安上がりになるのは
  $300 + n < 6n$ のとき，すなわちメッセージが 60 個を超えるときだけで
  ある．\tip{損益分岐点は，設定のコストを1回のやり取りあたりの節約で
  割った値である．ここでは $300/(6-1) = 60$ 個のメッセージとなる．}
\end{example}

Unix 系システムの IPC 機構とその利用は第9章の主題である．

\begin{keypoint}[章のまとめ]
  プロセスとは実行中のプログラムであり，そのアドレス空間---テキスト，
  データ，ヒープ，スタックからなり，ページテーブルによって仮想アドレス
  から物理アドレスへ変換される---と，オペレーティングシステムが PCB に
  保持する実行状態とによって表される．コンテキストスイッチは，ある
  プロセスの状態をその PCB へ退避し，別のプロセスの状態を復元するもので，
  直接のコストと，コールドキャッシュという間接のコストを伴う．プロセスは
  新規，実行可能，実行，待ち，終了の状態の間を移り，\texttt{fork} と
  \texttt{exec} によって生成される．CPU スケジューラは次に実行する
  実行可能なプロセスを選び，そのオーバーヘッドはタイムスライスに依存
  する．そして IPC は，メッセージパッシングまたは共有メモリによって，
  隔離されたプロセスの間の通信を可能にする．
\end{keypoint}

\end{document}
